% ===================================================================
% skeleton.tex -- Manuscript skeleton: thickness-dependent W-doped In2O3 (IWO) thin-film transistors studied by
% Silvaco ATLAS 5.28 device simulation calibrated on the fabrication group's data, Quantum ESPRESSO calculations of
% crystalline and amorphous slabs, and a pre-registered hypothesis register.
% Draft 2026-10-13 (Claude). Status: SKELETON FOR HUMAN REVIEW. Under charter section 9 no AI agent may put claims in a
% manuscript without explicit approval; nothing here is approved. Scope: a computational paper; claude/NATURE_PLAN.md
% stays a proposal (D-09) and no new measurements exist (D-11).
% Build: pdflatex skeleton; bibtex skeleton; pdflatex skeleton; pdflatex skeleton
% Placeholders: \pdPending{Pnn}{...}; every Pnn is listed in TODO.md with what it waits for.
% Source comments: % <!-- KEY:line -->; KEY legends are at the top of each input file. Extra keys used here:
%   PRIV = iwo-tft-research-private/   PKG = IWO_PHYSICS_CONSTRAINED_MODEL_V1/
%   CS = PRIV/01_CURRENT_STATE.md   LEDGER = PRIV/03_CLAIMS_LEDGER.yaml   DLOG = PRIV/05_DECISION_LOG.md
%   REG = PRIV/codex/work_chat/HYPOTHESIS_REGISTER.md   CHARTER = PRIV/00_PROJECT_CHARTER.md
%   NPLAN = PRIV/claude/NATURE_PLAN.md   PFIG = PKG/analysis_2026-10-06_extraction/PAPER_FIGURES.md
%   MEM = C:/Users/moham/.claude/projects/.../memory/iwo-paper-target.md (scope notes)
% ===================================================================
\documentclass[11pt]{article}
% ===================================================================
% paper_preamble.tex -- packages and macros for the paper draft (claude/paper_draft_2026-10-13).
% Copied from IWO_PHYSICS_CONSTRAINED_MODEL_V1/report_latex_full/preamble.tex (notation macros unchanged) and
% reduced to what an article needs. ASCII only. Writer-prefixed additions (\pd...) are at the end.
% Build: pdflatex skeleton; bibtex skeleton; pdflatex skeleton; pdflatex skeleton
% ===================================================================
\usepackage[T1]{fontenc}
\usepackage[utf8]{inputenc}
\usepackage[a4paper,top=2.5cm,bottom=2.5cm,left=2.5cm,right=2.5cm]{geometry}
\usepackage{setspace}
\onehalfspacing
\usepackage[expansion=false]{microtype}
\setlength{\emergencystretch}{3em}   % long typewriter tokens (tags, file names) in justified text
\usepackage{amsmath}
\usepackage{newtxtext,newtxmath}
\usepackage{siunitx}
\sisetup{separate-uncertainty=true, range-phrase={--}, range-units=single}
\usepackage{graphicx}
\usepackage{booktabs,longtable,array,tabularx}
\usepackage{caption}
\captionsetup{font=small,labelfont=bf,labelsep=period,format=plain}
\usepackage{enumitem}
\setlist{itemsep=2pt,topsep=4pt}
\usepackage[dvipsnames]{xcolor}
\usepackage[numbers,sort&compress]{natbib}
\usepackage[hidelinks]{hyperref}
\usepackage[capitalise,noabbrev]{cleveref}

% ---------------- summary / note environments (plain text, as in the report) ----------------
\newenvironment{keyresult}[1][]{\par\medskip\noindent\textit{Summary.}\ }{\par\medskip}
\newenvironment{caveat}[1][]{\par\medskip\noindent\textit{Limitation.}\ }{\par\medskip}

% ---------------- notation macros (identical to the report preamble) ----------------
\newcommand{\run}[1]{\texttt{run\_#1}}
\DeclareRobustCommand{\file}[1]{\texttt{\detokenize{#1}}}
\newcommand{\prov}[1]{\textsc{\MakeLowercase{#1}}}
\newcommand{\IWO}{IWO}
\newcommand{\InO}{In$_2$O$_3$}
\newcommand{\HfO}{HfO$_2$}
\newcommand{\AlO}{Al$_2$O$_3$}
\newcommand{\Vth}{V_{\mathrm{th}}}
\newcommand{\Vthcc}{V_{\mathrm{th,cc}}}
\newcommand{\Vg}{V_{\mathrm{G}}}
\newcommand{\Vd}{V_{\mathrm{D}}}
\newcommand{\Id}{I_{\mathrm{D}}}
\newcommand{\Ig}{I_{\mathrm{G}}}
\newcommand{\Ion}{I_{\mathrm{on}}}
\newcommand{\Ioff}{I_{\mathrm{off}}}
\newcommand{\gm}{g_{\mathrm{m}}}
\newcommand{\muFE}{\mu_{\mathrm{FE}}}
\newcommand{\muband}{\mu_{\mathrm{band}}}
\newcommand{\Cox}{C_{\mathrm{ox}}}
\newcommand{\Ec}{E_{\mathrm{c}}}
\newcommand{\EF}{E_{\mathrm{F}}}
\newcommand{\Eg}{E_{\mathrm{g}}}
\newcommand{\dEc}{\Delta E_{\mathrm{c}}}
\newcommand{\mstar}{m^{*}}
\newcommand{\Nc}{N_{\mathrm{c}}}
\newcommand{\Nt}{N_{\mathrm{t}}}
\newcommand{\NTA}{N_{\mathrm{TA}}}
\newcommand{\WTA}{W_{\mathrm{TA}}}
\newcommand{\Dit}{D_{\mathrm{it}}}
\newcommand{\Qf}{Q_{\mathrm{f}}}
\newcommand{\Nd}{N_{\mathrm{d}}}
\newcommand{\Dsr}{\Delta_{\mathrm{sr}}}
\newcommand{\kT}{k_{\mathrm{B}}T}
\newcommand{\mVdec}{\ensuremath{\,\mathrm{mV/dec}}}
\newcommand{\Aum}{\ensuremath{\,\mathrm{A}/\mu\mathrm{m}}}
\newcommand{\cmsq}{\ensuremath{\,\mathrm{cm}^{-2}}}
\newcommand{\cmcu}{\ensuremath{\,\mathrm{cm}^{-3}}}
\newcommand{\cmVs}{\ensuremath{\,\mathrm{cm}^{2}/\mathrm{V\,s}}}

% ---------------- writer-prefixed additions (paper draft only) ----------------
% \pdPending{Pnn}{KIND: text} marks a gap that waits for a result, a decision, a check or a citation.
% Every Pnn is listed in TODO.md with what it waits for. KIND is one of RESULT PENDING, DECISION PENDING,
% CHECK PENDING, CITATION PENDING, TEXT PENDING, DISCLOSURE PENDING.
\newcommand{\pdPending}[2]{\textbf{[#1~#2]}}
% \pdCell{name}: a DFT cell or run label in typewriter (underscores allowed).
\DeclareRobustCommand{\pdCell}[1]{\texttt{\detokenize{#1}}}
% \pdIntent{...}: the one-paragraph intent of a manuscript section (skeleton only).
\newenvironment{pdIntent}{\par\noindent\textit{Intent.}\ }{\par\medskip}


% Working title (not approved): "Registered device-simulation and first-principles tests of thickness-dependent
% W-doped In2O3 transistors". Title and abstract are written last, after the scorecard exists.
\title{\pdPending{P32}{TITLE PENDING: written after the scorecard (P30)}}
\author{\pdPending{P34}{DECISION PENDING: authors and affiliations}}
\date{Draft of 13 October 2026 -- not for circulation}

\begin{document}
\maketitle

\begin{abstract}
\pdPending{P33}{TEXT PENDING: abstract, written last}

\begin{pdIntent}
In about 150 words: the question (what sets the thickness dependence of ultrathin IWO transistors, and can a model with
shared parameters and computed inputs account for it); the approach (device simulation with one shared parameter set,
first-principles calculations of crystalline and amorphous films, all tests registered before their outcome); the
outcomes, including failed and unresolved predictions, stated with their labels; and the limits (one device per
thickness, retrospective device comparisons, PBE without $+U$, machine-learned structures that failed their force
check). It must not use the words ``predict'' or ``validate'' for anything fitted or retrospective.
\end{pdIntent}
\end{abstract}

\section*{Status of this draft}
This file is a skeleton. It fixes the section order and what each section is for, and it contains the Methods
(\cref{sec:methods}) and Limitations (\cref{sec:limitations}) as drafts. No claim in it is approved for a manuscript
(charter section 9). % <!-- CHARTER:105-116 -->
The paper is computational: it uses the four transfer curves already published by the fabrication group and adds no
measurement (decision D-11); the proposal for new devices (\file{claude/NATURE_PLAN.md}) is not part of it (decision
D-09). % <!-- DLOG:22 ; DLOG:20 ; NPLAN:3 -->
Results that are pending, sealed or blind appear only as placeholders.
\pdPending{P36}{DECISION PENDING: target journal; the charter treats a Nature-family paper as the ambition, not the
optimisation target, and the device data are already published}. % <!-- CHARTER:11-13 ; MEM:13-19 -->

% =====================================================================================================================
\section{Introduction}
\label{sec:intro}
\begin{pdIntent}
State why the thickness limit of oxide-semiconductor channels matters and what is known for IWO: the published
2--\SI{13}{nm} series with an analytic model of eight fitted parameters per device and crystalline DFT
\citep{Januar2026}. % <!-- CS:41-50 -->
Name the gap this paper addresses: whether a device model with one shared parameter set and inputs computed from first
principles can account for the thickness dependence, and what amorphous-film calculations say about W doping at the
thinnest thicknesses. Explain why every test was registered before its outcome, and say plainly that all device
comparisons are retrospective (four devices, all seen). End with the list of hypotheses (\cref{tab:m-reg-hyp}). Do not
present re-analysis of the published curves as the headline. % <!-- REG:37 ; MEM:19 -->
\end{pdIntent}

% =====================================================================================================================
\section{Results}
\label{sec:results}

\subsection{The measured devices under one extraction protocol}
\label{sec:r-data}
\begin{pdIntent}
Present the four transfer curves (one device each, $\Vd = \SI{0.7}{V}$) and the registered extraction: thresholds
0.662, 0.644 and \SI{0.083}{V} for 2.0, 6.3 and \SI{13.2}{nm} (claim C002, \prov{DERIVED}); off-state floors
(C003); the collapse of the local swing against $\Id/\Ion$ in measurement and simulation (C007, \prov{OBSERVED}); and
how the source paper's mobility and swing relate to our definitions (C030--C033). The \SI{31.8}{nm} device is shown as
context and model failure only. % <!-- LEDGER:11-20 ; LEDGER:40-44 ; LEDGER:178-205 ; LEDGER:172-176 -->
\end{pdIntent}
\pdPending{P35}{TEXT PENDING: figure plan for the main text and Extended Data (figure-maker), consistent with
PAPER\_FIGURES.md}. % <!-- PFIG:5-36 -->

\subsection{Calibrated models and the failed held-out test}
\label{sec:r-history}
\begin{pdIntent}
Summarise in one paragraph why the earlier per-film models are not predictive: thresholds and on-currents agree by
construction (C013, \prov{FITTED}), and every held-out \SI{6.3}{nm} threshold test failed, including the registered
V2 test (C014, \prov{FALSIFIED}). This motivates the shared-parameter campaign. % <!-- LEDGER:74-84 -->
\end{pdIntent}

\subsection{What sets the off-state floor}
\label{sec:r-floor}
\begin{pdIntent}
Report H-A, H-B and H-C as registered: strip ohmicity and transfer of the strip conductance from \SI{13.2}{nm} to the
thinner films; nominal and worst-case gate leakage; generation bounds; the S2-11 convergence check. Every verdict with
its censoring status (pass, fail, below resolution, provisional, unresolved). State what only a drain-bias measurement
could decide.
\end{pdIntent}
\pdPending{P05}{RESULT PENDING: S2-14..S2-17, STRIP-6p3} \pdPending{P06}{RESULT PENDING: S2-11 (sealed), S2-01..S2-13}

\subsection{Confinement, the threshold and the \texorpdfstring{\SI{6.3}{nm}}{6.3 nm} shape residual}
\label{sec:r-confinement}
\begin{pdIntent}
Report H-E (all confinement inputs off; affinity-only rigidity), including the failed S1-05 expectation and both
\SI{2.0}{nm} bands of A13, and H-F (two-film reference shift and scale, branch tests). Separate model-consistency
checks from statements about the device. % <!-- REG:184-212 ; REG:243-247 -->
\end{pdIntent}
\pdPending{P08}{RESULT PENDING: S3-01..S3-03, AFF-2p0, AFF-6p3} \pdPending{P09}{RESULT PENDING: S3-10..S3-13, S3-27; H-F scoring not yet frozen}

\subsection{One shared model across thickness}
\label{sec:r-shared}
\begin{pdIntent}
Report the S4 calibration as calibration, the identifiability analysis (which parameter combinations are determined),
the retrospective leave-one-film-out folds of H-J and the mobility scattering branch of H-D. If S4 reports the shared
model as failed, show the shift/shape decomposition and stop there: no per-film rescue parameters. % <!-- CD:324-338 ; REG:286-302 ; REG:164-177 -->
\end{pdIntent}
\pdPending{P07}{RESULT PENDING: S4, native triples} \pdPending{P10}{RESULT PENDING: S4a, S4c, S5}

\subsection{Crystalline first-principles inputs}
\label{sec:r-crystal}
\begin{pdIntent}
Give the computed confinement inputs and their use: gap opening of \SI{0.335}{eV} and mass $0.206\,m_0$ for the
\SI{2}{nm} slab, \SI{0.90}{eV} (\SI{0.95}{}--\SI{1.00}{eV} HSE-corrected) at \SI{1}{nm} (C020, \prov{COMPUTED}); the
conduction-band share of the opening; the W donor level in crystalline cells (C022); the sub-nanometre slabs as
preliminary (C021); and the oxygen-vacancy test H-L. State that these are crystalline while the films are amorphous or
nanocrystalline. % <!-- LEDGER:114-130 ; CS:211 -->
\end{pdIntent}
\pdPending{P16}{RESULT PENDING: H-L energies, k3/k4 check, A20 sensitivities}

\subsection{Amorphous films: W as an extended-carrier donor against thickness}
\label{sec:r-hn}
\begin{pdIntent}
Report H-G at \SI{2.0}{nm} (registered) and H-N (exploratory) as per-cell classes with every unresolved check, the
smearing family and lineage of each cell, and the linked disclosure footnote for the \SI{2.0}{nm} cells (A36 erratum).
State that the \SI{2.0}{nm} H-N endpoint is \prov{UNRESOLVED}, so neither a crossover nor its absence can be claimed,
and that a missing or exploratory \SI{0.8}{nm} cell also blocks a no-crossover statement. % <!-- REG:1016-1018 ; REG:1027 ; REG:900 -->
\end{pdIntent}
\pdPending{P23}{RESULT PENDING: H-N per-cell classes and the series outcome; H-G, H-H, H-K scores; release decision D-42
(not yet in the decision log)}

\subsection{Surfaces, gate-oxide interfaces and amorphous confinement}
\label{sec:r-interfaces}
\begin{pdIntent}
Report H-H (termination; \SI{2.0}{nm} part unresolved because the bare control was parked), H-I (expected
unresolved), H-M (alternative-interface control only) and H-K (amorphous against crystalline confinement), with the
gap-anchoring scheme of A23. % <!-- DLOG:46 ; REG:823-827 ; REG:396-397 ; REG:304-312 ; REG:539-544 -->
\end{pdIntent}
\pdPending{P24}{RESULT PENDING: A23 hybrid checks; amorphous bulk references}

\subsection{Scorecard of the registered predictions}
\label{sec:r-scorecard}
\begin{pdIntent}
One table with every registered prediction (H-A to H-M), its frozen criterion, the outcome and its status, and the
exploratory H-N read-out separately; failures and unresolved outcomes are listed in the same table as passes. % <!-- DLOG:28 -->
\end{pdIntent}
\pdPending{P30}{RESULT PENDING: the scorecard of every registered prediction and the H-N exploratory read-out}

% =====================================================================================================================
\section{Discussion}
\label{sec:discussion}
\begin{pdIntent}
Say what the registered outcomes establish and what they do not: which mechanisms are excluded within the model, which
remain open for lack of measurements (gate current, drain-bias series, device statistics), and which conclusions rest
on PBE alone. Distinguish calibrated, computed and tested statements throughout. List the measurements that would
decide the open questions (drain-bias dependence of the floor, devices per thickness, optical gap and Hall data),
without presenting them as planned work of this paper. % <!-- CS:217-220 ; DLOG:22 -->
\end{pdIntent}

% =====================================================================================================================
\section{Methods}
\label{sec:methods}
% ===================================================================
% methods_tcad.tex -- Methods: device simulation (Silvaco ATLAS 5.28), measured data, extraction protocol,
% campaign design and pre-registration, scoring/resolution method, execution controls, launch accounting.
% Draft 2026-10-13 (Claude), for human review; no claim here is approved for a manuscript (charter section 9).
% ASCII only. Every number carries a source comment  % <!-- KEY:line -->  (file:line). KEY legend:
%   PKG    = IWO_PHYSICS_CONSTRAINED_MODEL_V1/              PRIV   = iwo-tft-research-private/
%   GEOM   = PKG/config/geometry.yaml                       SOLV   = PKG/config/solver.yaml
%   PHYS   = PKG/config/physics_options.yaml                V3     = PKG/config/iwo_material_model_v3.yaml
%   PROT   = PKG/analysis_2026-10-06_extraction/PROTOCOL.md A1R    = PKG/analysis_2026-10-06_extraction/A1_RESULTS.md
%   PFIG   = PKG/analysis_2026-10-06_extraction/PAPER_FIGURES.md
%   CD     = PKG/analysis_2026-10-10_campaign/CAMPAIGN_DESIGN.md
%   S2PREP = PKG/analysis_2026-10-10_campaign/S2_PREP.md
%   S2FIX  = PKG/analysis_2026-10-10_campaign/S2_SCORER_FIXES_2026-10-12.md
%   S2GATE = PKG/analysis_2026-10-10_campaign/S2_RUNNER_GATES_2026-10-13.md
%   DEVTEX = PKG/report_latex_full/chapters/03_device_data.tex   V2TEX = PKG/report_latex_full/chapters/14_model_v2.tex
%   MAIN   = PKG/report_latex_full/main.tex
%   REG    = PRIV/codex/work_chat/HYPOTHESIS_REGISTER.md    LEDGER = PRIV/03_CLAIMS_LEDGER.yaml
%   DLOG   = PRIV/05_DECISION_LOG.md                        CS     = PRIV/01_CURRENT_STATE.md
% Not opened: PKG/results/runs/run_0075_* and PKG/results/sealed/ (S2-11 is sealed).
% ===================================================================
\subsection{Device simulation}
\label{sec:m-tcad}

\subsubsection{Devices and measured data}
\label{sec:m-tcad-data}

The simulated transistors are the bottom-gate W-doped \InO{} (IWO) thin-film transistors of the thickness series
published by \citet[pp.~2--3]{Januar2026}; the page is the one verified for the stack in
\file{report_latex_full/chapters/03_device_data.tex}. % <!-- DEVTEX:18-19 -->
The devices were made and measured by the fabrication group (National Taiwan University); this work adds no
measurement (decision D-11). % <!-- DLOG:22 ; CS:27 -->
From the gate upwards the stack is a \SI{50}{nm} TiN gate on n$^+$ Si, \SI{15.0}{nm} \HfO{}, \SI{2.0}{nm} \AlO{}, the
sputtered IWO channel (about 2\,\% W, composition basis not stated by the source) and \SI{70}{nm} Pd source and drain
contacts on top of the channel, with no passivation layer. % <!-- GEOM:4-9 ; DEVTEX:23-31 ; V3:19-21 -->
Both oxides were grown by atomic-layer deposition at \SI{250}{\celsius}, and the process ends with a rapid thermal
anneal at \SI{150}{\celsius} for \SI{600}{s}; these enter no parameter. % <!-- DEVTEX:25 ; DEVTEX:33-34 -->
\Cref{tab:m-tcad-stack} lists every structural item with its provenance label.

\begin{table}[htbp]\centering\footnotesize
\caption{Structural items of the simulated device and their provenance labels.}
\label{tab:m-tcad-stack}
\begin{tabularx}{\linewidth}{@{}l>{\raggedright\arraybackslash}Xl@{}}
\toprule
Item & Value & Label \\
\midrule
Gate & TiN \SI{50}{nm} (Si support screened, omitted) & \prov{MEASURED} (schematic) \\ % <!-- GEOM:3-4 -->
Gate dielectric & \HfO{} \SI{15.0}{nm}, $\varepsilon_r = 19.57$ & thickness from schematic; $\varepsilon_r$ \prov{MEASURED} (C--V, SI Fig.~S1) \\ % <!-- GEOM:5 -->
Interlayer & \AlO{} \SI{2.0}{nm}, $\varepsilon_r = 9.0$, touching the channel & thickness from schematic; $\varepsilon_r$ and position \prov{ASSUMED} \\ % <!-- GEOM:6 ; CD:168 -->
Channel & IWO, $t = 2.0$, 6.3, 13.2, \SI{31.8}{nm} (nominal), $\varepsilon_r = 9.3$ & \prov{MEASURED} ($\varepsilon_r$ from C--V, SI Fig.~S1) \\ % <!-- GEOM:7 ; V3:43-46 -->
Source/drain & Pd \SI{70}{nm}, top contacts, ideal Ohmic & geometry from schematic; contact model \prov{ASSUMED} \\ % <!-- GEOM:8 ; V3:209-211 -->
Contact overlap & \SI{2.0}{\micro\metre} per side & \prov{ASSUMED} \\ % <!-- GEOM:14 -->
$L$, $W$ & \SI{20}{\micro\metre}, \SI{290}{\micro\metre} & schematic \\ % <!-- GEOM:11-12 -->
Geometric $\Cox$ & \SI{8.9554e-7}{\farad\per\centi\metre\squared} & \prov{DERIVED} from the rows above \\ % <!-- PROT:20 -->
Gate work function & \SI{4.7}{eV} & \prov{ASSUMED} (TiN 4.5--4.9~eV); degenerate with $\Qf$ and $\chi$ \\ % <!-- V3:206-208 -->
Temperature & \SI{300}{K} & \prov{ASSUMED} (not recorded) \\ % <!-- V3:22-25 -->
Above the channel & air (no cap) & as specified \\ % <!-- GEOM:9 -->
\bottomrule
\end{tabularx}
\end{table}

\paragraph{Measured data.} One transfer curve per film was available: $\Vg$ from $-3$ to $+3$~V in \SI{50}{mV} steps
(121 points) at $\Vd = \SI{0.7}{V}$, for 2.0, 6.3, 13.2 and \SI{31.8}{nm}
(\prov{MEASURED}, claim C001). % <!-- CS:73-76 ; LEDGER:5-9 -->
There is no reverse sweep, no gate current, no second drain bias and only one device per thickness. % <!-- LEDGER:9 ; CS:77 -->
The thicknesses are nominal and the method of their measurement is unknown. % <!-- CS:253 -->
The \SI{31.8}{nm} device is not modelled (decision D-01); it never turns off ($\Id(-3\,\mathrm{V}) =
\SI{3.0e-7}{\Aum}$, claim C029, \prov{DERIVED}) and, having been in the data file since September, it is not a
holdout. % <!-- DLOG:12 ; LEDGER:172-176 -->
Every comparison between simulation and measurement in this work is therefore retrospective. % <!-- REG:37 ; CD:501 -->
The off-state floors, the median current over $\Vg = -2$ to $-0.5$~V, are $4.61\times10^{-15}$, $1.53\times10^{-13}$
and \SI{6.60e-12}{\Aum} for 2.0, 6.3 and \SI{13.2}{nm} (claim C003, \prov{DERIVED}). % <!-- LEDGER:16-20 ; CD:344 -->
\pdPending{P01}{DECISION PENDING: written permission of the fabrication group to use and publish the raw curves
(they are already in the public repository), data custody and authorship; charter section 8}.

\subsubsection{Simulation domain and numerical method}
\label{sec:m-tcad-numerics}

All device simulations use two-dimensional drift--diffusion in Silvaco ATLAS 5.28.1.R driven by DeckBuild 5.0.10.R
\citep{AtlasManual,DeckBuildManual}. % <!-- MAIN:35 -->
The domain contains the physical volumes: air above the channel, the Pd contacts, the \AlO{} and \HfO{} regions and
the TiN gate as a conductor volume. % <!-- GEOM:10 -->
The mesh width is \SI{1}{\micro\metre}, so currents are reported in A/\textmu m. % <!-- GEOM:13 -->
The $x$ mesh is refined to \SI{10}{nm} at the contact edges; along $y$ the spacing is \SI{0.0625}{nm} at the
IWO/\AlO{} interface, $\min(1.0\,\mathrm{nm}, t/8)$ in the film, \SI{0.25}{nm} in \AlO{} and \SI{1.0}{nm} in \HfO{}. % <!-- GEOM:18 ; GEOM:20-23 -->
The interface trap sheet is regularised as a \SI{0.25}{nm} IWO layer; halving that layer changes $\Id$ by
\SI{0.009}{dec}. % <!-- GEOM:19 -->

The solver uses Fermi--Dirac statistics, a Shockley--Read--Hall background and electrons only
(\texttt{carriers=1}); the hole current is numerically zero in every recorded run. % <!-- SOLV:2-3 -->
The Newton scheme includes trap occupancy (\texttt{newton trap}, \texttt{maxtraps=10}, \texttt{itlimit=80},
\texttt{climit=1e-6}, \texttt{ir.tol=1e-19}). % <!-- SOLV:4 -->
In depletion, convergence is accepted with \texttt{XANDRNORM} and a continuity tolerance
$\max(10^{-17}\,\mathrm{A},\ 10^{-5}\times$ the smallest measured active current$)$, up to a measured current of
\SI{1e-9}{\Aum}. % <!-- SOLV:5-8 -->
The trap density of states is integrated on 384/192 levels in the campaign (768/384 in the S1-04 check). % <!-- CD:195 -->
Gate sweeps are stepped \texttt{SOLVE VSTEP/VFINAL} blocks from $-3$~V: \SI{0.2}{V} steps to $-1.0$~V,
\SI{0.05}{V} to \SI{1.5}{V} and \SI{0.1}{V} to \SI{3.0}{V}, after a \SI{0.1}{V} ramp to the sweep start. % <!-- SOLV:10-13 -->
Launches run strictly one at a time (parallel launches failed), with thickness-dependent time limits of 1200, 1500,
2400 and \SI{3000}{s} for 2.0, 6.3, 13.2 and \SI{31.8}{nm} (multiplied per entry where listed). % <!-- CS:285 ; CD:289 ; CD:479 -->
A run passes the current-balance (KCL) check when
$|I_D + I_S + I_G| \le \max(10^{-17}\,\mathrm{A},\ 0.1\times\mathrm{floor}) + 1\,\%\ \sum|I|$. % <!-- CD:356 ; SOLV:19-21 -->
Three tool facts constrain the decks: \texttt{MOBILITY} statements placed between \texttt{SOLVE}s are ignored, at most
ten trap distributions are allowed, and \texttt{DEFECTS} and \texttt{INTDEFECTS} take densities per volume and per
area, respectively. % <!-- CS:288-290 -->

\paragraph{Numerical verification.} Before any hypothesis-scoring launch, the baseline physics was checked under
refinement: mesh $\times0.7$ at \SI{13.2}{nm} (S1-03) changed the active log-current error by $-0.0005$~dec and a
doubled trap-DOS grid at \SI{2.0}{nm} (S1-04) by $-0.00002$~dec, both inside the pre-registered numerics rule
(\SI{0.01}{dec}, \SI{5}{mV} in $\Vthcc$, \SI{3}{mV/dec} in SS). % <!-- REG:29-30 ; CD:298 -->
A native launch (S1-07) confirmed the offline shift/scale transformation used for profiling to a maximum deviation of
$2.05\times10^{-4}$~dec; this validates that transformation only. % <!-- REG:33 -->
The extraction carries a separate numerical budget of $\pm1.5$~mV in $\Vth$, $\pm1$~mV/dec in SS and $\pm0.7\,\%$ in
$\gm$. % <!-- PROT:67 -->

\paragraph{TCAD check rule.} Every change of the deck generator or of a scoring rule is reviewed by Codex and checked
by hand against the ATLAS 5.28 manual, which governs, with the 2004 manual as a second reference; manual pages are
quoted where a model's numerical control is at stake (for example the oxide mesh as the accuracy control of
self-consistent tunnelling). % <!-- S2PREP:158-160 ; S2PREP:174 -->
A regression script rebuilds every launched deck from its recorded overrides and must report no unexpected
difference before the next launch. % <!-- CD:290 -->

\subsubsection{Physical model}
\label{sec:m-tcad-model}

The campaign model (V3) uses one parameter set for the 2.0, 6.3 and \SI{13.2}{nm} films; film dependence enters only
through the thickness $t$ in physical laws. % <!-- CD:24-27 ; CD:487-489 -->
Its ingredients and their labels are listed in \cref{tab:m-tcad-params}.
\begin{itemize}
\item \emph{Confinement.} The conduction-band shift $\dEc(t)$ (in \texttt{AFFINITY}), the gap opening
$\Delta E_g(t)$ (in \texttt{EG300}) and the mass increment $\Delta m^*(t)$ (in \texttt{NC300}) follow power laws fitted
to our crystalline PBE slab points at H--H thicknesses 0.695, 0.930 and \SI{1.962}{nm}; no device data enter these
laws (\prov{FITTED} to \prov{COMPUTED} points). % <!-- V3:58-74 ; V3:84-97 ; V3:154-172 -->
Values at 6.3 and \SI{13.2}{nm} are extrapolations, and identifying the device thickness with the slab H--H thickness
is \prov{ASSUMED}. % <!-- CD:151 -->
The \SI{0.695}{nm} point comes from a slab whose result is still marked preliminary and is not yet entered in the DFT
results log \pdPending{P02}{CHECK PENDING: enter the 0.75 nm (H--H 0.695 nm) slab in RESULTS\_LOG.md, or refit the
laws without it; claim C021}. % <!-- LEDGER:120-124 ; V3:68 -->
\item \emph{Quasi-two-dimensional density of states.} At \SI{2.0}{nm} only, $\Nc$ is replaced by $N_{2D}/t$
(\prov{CALCULATED}) under a pre-registered single-subband rule (DFT subband spacing \SI{0.7191}{eV} $\gg \kT$). This
substitution inside the 3-D Fermi--Dirac form overstates the degenerate on-state density by up to $1.7\times$. % <!-- V3:98-105 ; CD:63 ; CD:111 -->
\item \emph{Traps.} An exponential acceptor tail with integrated density $N_s$ spread uniformly over the film
($\Nt = N_s/t$, $\NTA = \Nt/\WTA$); a deep acceptor Gaussian; and an interface acceptor sheet at the IWO/\AlO{}
interface. Trap levels are referenced to $\Ec$. % <!-- V3:106-118 ; V3:178-195 -->
\item \emph{Charges and donors.} A uniform front fixed charge $\Qf$ at $y = 0$ carries the alignment offset
$\Delta V_0$; it is degenerate with the gate work function and the bulk affinity, which are therefore held fixed.
The donor density is constant: the fitted V1 amplitude is kept, and the constant form is \prov{ASSUMED}. % <!-- CD:50 ; CD:116 ; V3:119-130 -->
\item \emph{Mobility.} A constant band mobility per film from one shared Matthiessen law,
$\muband(t) = \mu_\infty\,(0.208\,m_0/\mstar(t))/(1+\ell/t)$, multiplied by the roughness factor $(1-\Dsr/t)^2$ with
$\Dsr = \SI{0.287}{nm}$ from the source paper. % <!-- V3:131-153 -->
\end{itemize}
Models switched off, each with a stated reason, are field-dependent mobility, the silicon low-field and inversion-layer
mobility laws, band-to-band tunnelling, Coulombic trap barrier lowering, lumped contact resistance, a back-surface
charge sheet and the density-gradient model. % <!-- CD:129-137 -->
Branches tested once and adopted only by a pre-registered rule are the IGZO.TOKYO carrier-density mobility, Schottky
Pd contacts with a shared barrier from $\{0.2, 0.5\}$~eV, and a Bohm-potential pair test at \SI{2.0}{nm}. % <!-- CD:118-120 ; CD:308-321 -->

\begin{table}[htbp]\centering\footnotesize
\caption{Campaign model V3. Values at 2.0 / 6.3 / \SI{13.2}{nm}. Fitted values are the S4 start values; they are not
calibration results. Labels follow the claims-ledger vocabulary.}
\label{tab:m-tcad-params}
\begin{tabularx}{\linewidth}{@{}>{\raggedright\arraybackslash}p{3.3cm}>{\raggedright\arraybackslash}X>{\raggedright\arraybackslash}p{4.6cm}@{}}
\toprule
Parameter & Value & Label and basis \\
\midrule
Bulk gap, affinity & 3.05~eV; 4.3~eV & \prov{LITERATURE} (another sputtered IWO device) \\ % <!-- V3:27-37 -->
Bulk mass & $0.208\,m_0$ & \prov{LITERATURE} (optical Hall, \InO{}) \\ % <!-- V3:38-42 -->
$\dEc(t)$ & $0.70147\,t^{-1.31669}$~eV: 0.2816 / 0.0622 / 0.0235~eV & \prov{FITTED} to \prov{COMPUTED} DFT points \\ % <!-- V3:155-159 -->
$\Delta E_g(t)$ & $0.81431\,t^{-1.31518}$~eV & \prov{FITTED} to \prov{COMPUTED} DFT points \\ % <!-- V3:60-63 -->
$\mstar(t)$ & $0.208 + 0.11116\,t^{-1.26387}$: 0.2543 / 0.2189 / 0.2123~$m_0$ & \prov{LITERATURE} + \prov{FITTED} increment \\ % <!-- V3:84-89 -->
$\Nc$ (\SI{300}{K}) & $1.373\times10^{19}$ (quasi-2D) at \SI{2.0}{nm}; 3-D formula otherwise & \prov{CALCULATED} \\ % <!-- V3:98-105 -->
Tail $N_s$, $\WTA$ & $4\times10^{12}$\cmsq{}, \SI{40}{meV} (start) & \prov{FITTED} shared (S4) \\ % <!-- V3:106-116 -->
Deep acceptor Gaussian & $5\times10^{16}$~cm$^{-3}$\,eV$^{-1}$ at $\Ec - 0.6$~eV, width 0.15~eV & \prov{ASSUMED}; zero in S6 \\ % <!-- V3:178-182 -->
Interface acceptor sheet & $3\times10^{11}$~cm$^{-2}$\,eV$^{-1}$ at $\Ec - 0.3$~eV, width 0.12~eV & \prov{ASSUMED}; zero in S6 \\ % <!-- V3:189-195 -->
Donors $\Nd$ & $2.5\times10^{17}$\cmcu{}, constant & \prov{FITTED} (V1 amplitude, kept); form \prov{ASSUMED} \\ % <!-- V3:119-128 -->
Front charge $\Qf$ & $1.38345\times10^{12}$\cmsq{} (start) & \prov{FITTED} shared, profiled \\ % <!-- V3:196-199 -->
Mobility law & $\mu_\infty = 100.0$\cmVs{}, $\ell = 7.25$~nm (start) & \prov{FITTED} shared, profiled \\ % <!-- V3:147-153 ; CD:178-179 -->
Capture cross-sections, SRH lifetime & $10^{-15}$~cm$^2$; $10^{-6}$~s & \prov{ASSUMED} \\ % <!-- V3:176-177 ; V3:54-57 -->
Valence-band DOS, hole mobility & $10^{19}$\cmcu{}; 0.01\cmVs{} & \prov{ASSUMED} (holes not solved) \\ % <!-- V3:47-53 -->
\bottomrule
\end{tabularx}
\end{table}

\subsubsection{Model history: per-film calibration and the registered held-out test}
\label{sec:m-tcad-history}

Model V1 was calibrated per film: $\Qf$ to the constant-current threshold and the band mobility to $\Ion$, with the
tail width fitted on the \SI{2.0}{nm} subthreshold region (\prov{FITTED}). % <!-- CS:141-148 -->
Model V2 replaced the literature confinement laws by our DFT laws. Its pass/fail criteria were committed before the
launches (commit \texttt{7b75ee6}); one error in the registered transformation was corrected and published
(\texttt{a914165}) before the held-out outcome existed, and both versions were scored (\texttt{852dc8a}). % <!-- CS:164 ; V2TEX:43-50 -->
The held-out \SI{6.3}{nm} threshold prediction failed in both variants (predicted $\Vthcc$ 0.270 and 0.417~V against
0.644~V measured), as did the earlier shared-parameter V1 test (claim C014, \prov{FALSIFIED}). % <!-- LEDGER:80-84 ; CS:173-178 -->
The V2 off-state strip, a gate-independent parallel conductance with one geometric factor fitted on \SI{13.2}{nm}, is
a \prov{HYPOTHESIS} (claim C006). % <!-- LEDGER:34-38 ; V2TEX:137-145 -->
V1 and V2 enter this paper only as the documented starting point of the campaign below.

\subsubsection{Extraction protocol}
\label{sec:m-tcad-extraction}

Threshold, subthreshold swing and mobility are extracted by one code path for measured and simulated curves, under a
protocol written and committed (\texttt{7a61687}) before any measured-versus-TCAD number was computed. % <!-- PROT:3 ; CS:86 ; DLOG:16 -->
\begin{itemize}
\item \emph{Threshold (primary).} $\Vthcc$ is $\Vg$ at $\Id = \SI{1e-9}{\Aum}$, by log-linear interpolation at the
last upward crossing below +3~V; \SI{5e-10}{\Aum} is also reported. % <!-- PROT:32 -->
\item \emph{Subthreshold swing (registered window).} One decade between $5\times10^{-11}$ and \SI{5e-10}{\Aum},
admissible when the lower edge is at least ten times the floor; otherwise a floor-corrected value is primary, which
assumes an additive, gate-independent parallel current. % <!-- PROT:38-43 -->
\item \emph{Amendment A1 (secondary window).} Because the films differ in current scale, a window normalised by each
film's measured on-current, $[r, 10r]\times\Ion$ with $r = 10\times\max_t(\mathrm{floor}/\Ion) = 2.15\times10^{-5}$,
was added (commit \texttt{7ce7116}) after the measured values under the registered rule had been computed and before
any simulated value under the new rule existed; the disclosure is part of the amendment. % <!-- PROT:80-98 ; A1R:3 ; CS:87 -->
\item \emph{Mobility.} $\muFE = L\,\gm/(W\Cox\Vd)$ from central-difference $\gm$, reported at 1.0--3.0~V; it equals a
carrier mobility only where $\Vg - V_T \ge \Vd$ and the mobility is gate independent, and is otherwise a normalised
transconductance. A gradual-channel power-law fit $\mu = \mu_0 (u/1\,\mathrm{V})^\gamma$ over 1.0--3.0~V gives a
model-based channel mobility. % <!-- PROT:50-58 -->
\item \emph{Uncertainties and comparison.} Measurement noise from local quadratic residuals, interpolation-method
differences and the TCAD numerical budget are combined in quadrature; a difference is ``consistent'' if it is within
$2\sigma$. $\Cox$, the charge-centroid ratio $C_{\mathrm{eff}}/\Cox = 0.87$, contact resistance and device-to-device
spread are stated separately and not included. % <!-- PROT:63-78 -->
\end{itemize}
The source paper's peak mobilities are reproduced numerically on our curves by the peak of the saturation-formula
mobility (claim C030, \prov{DERIVED}); that this is the paper's method is only the leading candidate (C031,
\prov{INFERRED}), and the paper's near-thermal swing at \SI{2}{nm} remains \prov{UNRESOLVED} (C025). We report our
definitions throughout. % <!-- LEDGER:178-191 ; LEDGER:146-152 -->
The \SI{13.2}{nm} swing mismatch is reported as a finding; no window or floor treatment was chosen to remove it
(decision D-06). % <!-- DLOG:17 -->
\pdPending{P03}{CHECK PENDING: the fabrication group's mobility and SS recipe (formula, $\Cox$, $W/L$, current
normalisation, sweep provenance); claims C024, C025, C031}.

\subsubsection{Campaign design}
\label{sec:m-tcad-campaign}

The campaign (C10) was designed before its first launch from three independent inputs (the DFT-to-TCAD extraction, an
audit of the deck generator and ATLAS 5.28, and an independent design by Codex), merged, and reviewed read-only by
Codex. % <!-- CD:3-10 -->
Its budget was 200 sequential launches (decision D-13). % <!-- DLOG:23 ; CD:39-41 -->
\begin{itemize}
\item \emph{Shared fitted parameters.} Five values for all three films: $\Delta V_0$ (written as one $\Qf$), $\mu_\infty$,
$\ell$, $N_s$ and $\WTA$, plus a sixth only if a mobility branch is adopted. No per-film charge or mobility exists. % <!-- CD:24-27 -->
\item \emph{Exact profiling.} A uniform front charge shifts a transfer curve rigidly, and with a constant (or
IGZO-type) mobility the current is proportional to \texttt{MUN}; both hold in the continuum equations up to the
generation term, bounded at \SI{1e-28}{\Aum} against currents of at least \SI{1e-15}{\Aum} (\prov{CALCULATED}). % <!-- CD:228-233 -->
$\Delta V_0$, $\mu_\infty$ and $\ell$ are
therefore profiled offline on each launched film triple, without extrapolation beyond the solved $\Vg$ range, and every
optimum is confirmed natively within \SI{0.01}{dec}. % <!-- CD:28-32 ; CD:225-243 -->
\item \emph{Objective.} $J = \tfrac13\sum_t (\mathrm{RMSE}^2_{\mathrm{sub},t} + \mathrm{RMSE}^2_{\mathrm{on},t})/2$,
with log residuals over the measured subthreshold range from ten times the floor to \SI{1e-9}{\Aum} and over
$\Vg = 1$--3~V; no film is weighted by current or point count. % <!-- CD:215-223 -->
\item \emph{Acceptance} (\prov{ASSUMED} engineering limits, not confidence intervals). Per film: active log-RMSE
$\le \SI{0.15}{dec}$, $|\Delta\Vthcc| \le \SI{0.10}{V}$, $\Id(3\,\mathrm{V})$ within $\pm30\,\%$, SS in the A1 window
within $\pm20\,\%$, $\muFE$ at 1.5, 2 and 3~V within $\pm20\,\%$. Success on all three films is still called
calibration. % <!-- CD:246-255 -->
\item \emph{Identifiability.} A scaled residual Jacobian from the calibration archive; a condition number above 100 or
$|\mathrm{corr}| > 0.95$ (\prov{ASSUMED}) means only a parameter combination is reported. % <!-- CD:257-261 -->
\item \emph{Stages.} S0 smoke (2 launches), S1 baselines and numerics (7), S2 off-floor mechanisms (17), S3 model-form
branches (27), S4 calibration (45 space-filling + 18 adaptive + 3 native confirmation + 3 runner-up), S5
leave-one-film-out (30, retrospective), S6 sensitivities and frozen predictions (27), S7 \SI{31.8}{nm} stress (2),
reserve ($\le 19$). % <!-- CD:269-282 -->
\item \emph{Stop rules.} Among them: if the best $\sqrt J$ after S4a exceeds \SI{0.30}{dec} with systematic residuals,
the shared model is reported as failed and no per-film rescue parameter is introduced; a held-out film that fails the
acceptance falsifies transferability of the model form and is not repaired. % <!-- CD:324-338 -->
\end{itemize}

\paragraph{Off-state floor tests.} Gate leakage is computed in the solver with self-consistent tunnelling
(\texttt{QTNLSC.EL}), a second Fowler--Nordheim formulation (\texttt{FNORD}) and defect-assisted oxide tunnelling
(\texttt{ITAT.SC.EL}); generation is bounded with extreme SRH and trap-assisted tunnelling in two-carrier runs; and an
ungated IWO strip is simulated as a separate deck. % <!-- CD:34-38 ; CD:121-128 -->
Two oxide parameter sets are used: a nominal set from the literature (\HfO{}/\AlO{} affinities 2.5/1.25~eV, gaps
5.8/8.8~eV, tunnelling masses \prov{ASSUMED} inside literature ranges) and a deliberately leakage-maximising set
(TiN/\HfO{} barrier 1.5~eV, \prov{ASSUMED}), which is an optimistic scenario, not a rigorous upper bound. % <!-- CD:200-206 ; CD:82 ; CD:376 ; PHYS:58-78 -->
Because \texttt{FNORD} evaluates the field in the \AlO{} at the semiconductor interface, the \HfO{}-barrier constants
are transferred to that field ($r = \varepsilon_{\mathrm{HfO_2}}/\varepsilon_{\mathrm{Al_2O_3}} = 2.174$). % <!-- CD:81 ; PHYS:34-40 -->
\pdPending{P04}{CITATION PENDING: oxide band parameters (Robertson 2000, J. Vac. Sci. Technol. B 18, 1785, as cited in
the campaign design), tunnelling-mass sources, Lenzlinger--Snow constants, Ielmini trap-assisted tunnelling,
IGZO.TOKYO (Abe 2014); not yet in the bibliography}.
A mechanism explains a film's floor only if at least 80\,\% of the window points are valid, its own contribution
(against a matched control) lies within \SI{0.5}{dec} of the floor, its slope is within twice the measured slope plus
$2\sigma$, and the current balance passes; it explains the thickness scaling only if the three films agree within
\SI{0.5}{dec}. At \SI{2.0}{nm} an under-prediction does not falsify a mechanism, because that floor may be
instrument-limited. % <!-- CD:358-371 -->

\subsubsection{Pre-registered TCAD predictions}
\label{sec:m-tcad-register}

Before any S2 launch, the hypothesis register (\cref{sec:m-reg}) froze which launches score which hypothesis, the
predictions, the stop rules and the calibration membership. % <!-- REG:1089-1095 ; DLOG:28 -->
\Cref{tab:m-tcad-reg} summarises the registered TCAD predictions; every one is scored and reported whatever the
outcome. The diagnostic set comprises 26 launches, with the oxide-mesh check S2-11 first. % <!-- REG:1031-1045 ; REG:1054-1064 -->

\begin{table}[htbp]\centering\footnotesize
\caption{Registered TCAD predictions (register revision 2 and its amendments). Every threshold is an \prov{ASSUMED}
decision threshold. Outcomes are not part of this draft.}
\label{tab:m-tcad-reg}
\begin{tabularx}{\linewidth}{@{}l>{\raggedright\arraybackslash}X>{\raggedright\arraybackslash}p{3.4cm}@{}}
\toprule
ID & Registered prediction & Outcome \\
\midrule
H-B & Strip $I(0.7\,\mathrm{V})/I(0.1\,\mathrm{V}) = 7.0\pm0.5$; with $(W/L)_{\mathrm{par}}$ fitted on \SI{13.2}{nm} only, the floors at 6.3 and \SI{2.0}{nm} within \SI{0.5}{dec} of the measured ones & \pdPending{P05}{RESULT PENDING: S2-14..S2-17, STRIP-6p3} \\ % <!-- REG:126-131 -->
H-C & Nominal leakage at $-3$~V $\le 10^{-2}\times$ floor (only with a demonstrated ceiling below that bound); extreme generation $\le \SI{1e-20}{\Aum}$ or below resolution; worst-case leakage fails flatness or thickness scaling; resolved leakage changes by $\le\SI{0.1}{dec}$ under the S2-11 oxide-mesh refinement & \pdPending{P06}{RESULT PENDING: S2-11 (sealed), S2-01..S2-13} \\ % <!-- REG:139-151 -->
H-D & Band-mobility ratio \SI{2.0}{nm}/\SI{13.2}{nm}: mass-only 0.75--0.95 (reported ``unresolved: not computed'', A22); scattering branch 0.20--0.35 with $\ell = 4$--\SI{12}{nm} (crystalline mass inputs) & \pdPending{P07}{RESULT PENDING: S4, native triples} \\ % <!-- REG:164-176 -->
H-E & All confinement inputs off: $\Delta\Vthcc = -0.13\pm0.03$~V at \SI{2.0}{nm} (A13; the original $-0.23\pm0.04$~V is also scored), $-0.062\pm0.015$~V at 6.3, $-0.023\pm0.008$~V at \SI{13.2}{nm}; S3-01 $-$ S1-05 $= -0.272\pm0.010$~V; affinity-only shifts $-281.6\pm5$ and $-62.2\pm5$~mV, rigid within \SI{5}{mV} & \pdPending{P08}{RESULT PENDING: S3-01..S3-03, AFF-2p0, AFF-6p3} \\ % <!-- REG:184-204 -->
H-F & Two-film reference: \SI{6.3}{nm} needs a rigid shift of $+0.28\pm0.08$~V and a current scale of 0.15--0.35; baseline shape RMSE $> \SI{0.15}{dec}$ & \pdPending{P09}{RESULT PENDING: S3-10..S3-13, S3-27; H-F scoring not yet frozen} \\ % <!-- REG:243-247 ; S2PREP:167-173 -->
H-J & Every film: RMSE $\le\SI{0.15}{dec}$, $\Vth$ within $\pm0.10$~V, $\Ion$ within $\pm30\,\%$, SS and $\muFE$ within $\pm20\,\%$; any failure falsifies transfer & \pdPending{P10}{RESULT PENDING: S4a, S4c, S5} \\ % <!-- REG:286-293 -->
\bottomrule
\end{tabularx}
\end{table}

\subsubsection{Scoring and the resolution method}
\label{sec:m-tcad-scoring}

All registered TCAD scores come from one frozen scorer (\file{score_register.py}, sha256 \texttt{c0f28bb9}\ldots
\texttt{b55a}; self-test 90/90). % <!-- S2FIX:349 ; S2GATE:222 ; S2GATE:243 -->
It applies the register's censoring rule: a quantity below the demonstrated numerical resolution is reported as
\emph{below resolution}, counts as consistent with an upper bound only if that resolution is itself below the bound, and
is otherwise \emph{unresolved}, never a pass. % <!-- REG:16-20 -->
The register does not define ``demonstrated resolution''. The scorer therefore uses an operational method, which we
label \prov{ASSUMED}:
\begin{itemize}
\item a resolution is demonstrated only by a pair of two distinct completed runs that differ by one verified numerical
intervention (trap-DOS levels, global mesh factor, or oxide $y$-mesh factor), checked against the actual deck values
and recomputed from the run files, which are bound by sha256, at every use; % <!-- S2FIX:20 ; S2FIX:150 -->
\item the pair must agree at least at three points above the candidate resolution, one of them within one decade, over
the full bias range, with a deck-ratio tolerance of $10^{-3}$; % <!-- S2FIX:162-170 -->
\item a resolution applies only to its own quantity, film, physical settings, bias range and numerical endpoint; it is
never transferred between films, and for an upper-bound test the pair may not contain a tested run; % <!-- S2FIX:156-160 ; S2FIX:273 -->
\item a verdict computed with the runner's current-balance limit as a proxy is labelled \emph{provisional} and can
never pass an upper bound. % <!-- S2PREP:122 ; S2FIX:23 -->
\end{itemize}
Pair agreement establishes numerical stability under the tested intervention, not an absolute error bound. The
one-decade support value was chosen knowing the two numerical pairs already seen in S1; it is prospective for S2 only. % <!-- S2FIX:288-291 ; S2FIX:167 -->
With the current run plan, almost every primary S2 and affinity-only verdict will therefore be provisional or
unresolved. % <!-- S2FIX:238-247 -->
\pdPending{P11}{DECISION PENDING: the user's sign-off of the resolution method as a logged amendment, required before
the S2-11 outputs are unsealed}. % <!-- S2FIX:129 ; S2GATE:248 -->
\pdPending{P12}{DECISION PENDING: whether to add matched numerical refinement runs (a logged plan change) so that
primary S2 verdicts can become demonstrated}. % <!-- S2FIX:247 -->
The \SI{6.3}{nm} shape-residual hypothesis (H-F) has no frozen operational scoring yet; its five launches are held. % <!-- S2PREP:137 ; S2PREP:167-173 -->
The oxide trap-assisted tunnelling runs (S2-12, S2-13) cannot count as converged without a further logged amendment
(for example a \texttt{TAT.SLICEPTS} sensitivity launch), because the S2-11 verdict covers \texttt{QTNLSC} only. % <!-- S2PREP:164 -->

\subsubsection{Execution controls, sealing and protocol deviations}
\label{sec:m-tcad-controls}

The campaign runner enforces the budget cap, sequential execution and refusal of held plans in code. % <!-- CD:288-289 ; S2PREP:37 -->
For the diagnostic stage, executable stop gates evaluate the frozen leakage and generation scorers in an isolated
process and expose only a state (fired, not fired, unresolved); every gate decision is bound by sha256 to its
declaration, the scorer, the scored run files and the cited evidence, and an unresolved or fired gate stops the plan
until an audit file citing that decision is supplied. % <!-- S2GATE:32-60 ; S2GATE:186-207 -->
A state-only mode keeps metrics, error texts and launch output in a sealed side record, so that run status can be
followed without exposing results; timeout clean-up terminates only identified processes of the launch and never the
licence daemons. % <!-- S2GATE:112-123 ; S2GATE:82-88 ; S2GATE:186-191 -->
\pdPending{P13}{CHECK PENDING: the runner revision 3 (stop gates, sealing, clean-up) was uncommitted and DIAG not
launched as of S2\_RUNNER\_GATES\_2026-10-13.md; record the commit and the Codex verdict}. % <!-- S2GATE:3 -->

The following deviations are disclosed.
\begin{itemize}
\item The smoke run S0-02 needed six launches against a cap of three repairs. % <!-- REG:39 -->
\item The regression gate read five unexpected differences while runs \run{0067} to \run{0074} were launched, against
the rule that it must read zero; the five were reclassified as known after line-by-line verification on 2026-10-11
(review item R23). % <!-- S2PREP:82 -->
\item The registered S1-05 expectation ($+40$ to $+70$~mV) failed ($+142$~mV) and stays recorded as a failure; the
\SI{2.0}{nm} confinement band was re-registered after S1 had been seen and before any S3 run, and both bands are scored
(A13). % <!-- REG:31 ; REG:184-190 ; REG:211-212 -->
\item The two-carrier check S1-06 timed out and is unresolved. % <!-- REG:32 -->
\item S2-11 (\run{0075}) has been launched; its outputs are sealed and were not opened for this draft. Its eligibility
under the timeout rule and its ``possibly exposed'' label await rulings. % <!-- S2GATE:11-16 ; S2GATE:123 ; S2GATE:212 ; S2GATE:249 -->
\pdPending{P14}{DECISION PENDING: run\_0075 rulings (timeout eligibility, ``possibly exposed'' label) and the
unsealing of S2-11}.
\end{itemize}

\subsubsection{Launch accounting}
\label{sec:m-tcad-launches}

Every launch counts, failures included. % <!-- CD:287 -->
\Cref{tab:m-tcad-launches} gives the accounting as recorded.

\begin{table}[htbp]\centering\footnotesize
\caption{ATLAS launch accounting (global counter and the C10 campaign counter).}
\label{tab:m-tcad-launches}
\begin{tabularx}{\linewidth}{@{}>{\raggedright\arraybackslash}p{4.6cm}>{\raggedright\arraybackslash}X@{}}
\toprule
Stage & Launches \\
\midrule
V1 report (2026-09-25) & 52 recorded launches \\ % <!-- MAIN:13 ; MAIN:35 -->
Before the campaign (2026-10-07) & 60 of 71 approved \\ % <!-- CS:1 ; CS:276 -->
Campaign approval (D-13) & 200 more; global cap 260 $= 60 + 200$, enforced in code \\ % <!-- DLOG:23 ; SOLV:18 -->
At the register freeze & 14 C10 launches used (\run{0061}--\run{0074}), 186 remaining \\ % <!-- REG:38 -->
As of 2026-10-13 & 75 of 260 global; 15 of 200 C10 (including S2-11, \run{0075}) \\ % <!-- S2GATE:1 ; S2GATE:9 -->
Planned diagnostic set & 26 launches; C10 use 40 of 200 at its end \\ % <!-- REG:1038 ; S2PREP:141 -->
Projected total & about 110--140 launches \\ % <!-- REG:1066 -->
Final count & \pdPending{P15}{RESULT PENDING: final launch count and status of every launch} \\
\bottomrule
\end{tabularx}
\end{table}

% ===================================================================
% methods_dft.tex -- Methods: first-principles calculations (Quantum ESPRESSO): crystalline references, the amorphous
% workflow (MACE melt-quench, slab construction, QE settings, smearing families and amendments A15/A26/A31/A35/A36,
% D-24 restart/guard rules, export and certification), read-out definitions and CERN resources.
% Draft 2026-10-13 (Claude), for human review; no claim here is approved for a manuscript (charter section 9).
% SIGHT RULE: no amorphous-slab electronic quantity (classes, PR_n, gaps, edges, Fermi levels, projections, band energies)
% appears here. Only what RESULTS_LOG.md, the claims ledger and the register text already log was used; FINAL exports,
% XML/projection outputs, relaxation outputs and scratchpad quarantine folders were NOT opened.
% ASCII only. Every number carries a source comment  % <!-- KEY:line -->  (file:line). KEY legend:
%   PKG    = IWO_PHYSICS_CONSTRAINED_MODEL_V1/              PRIV   = iwo-tft-research-private/
%   DFTTEX = PKG/report_latex_full/chapters/13_dft.tex       RLOG   = PKG/qe_workflow/cern_htcondor/results/RESULTS_LOG.md
%   V3     = PKG/config/iwo_material_model_v3.yaml
%   SLAB   = PRIV/compute/amorph/SLAB_PROTOCOL.md            MLV    = PRIV/claude/plans/ML_VALIDATION.md
%   MACE   = PRIV/compute/amorph/mace_amorph.py              MKSPEC = PRIV/compute/amorph/make_specs.py (docstring)
%   MKDFT  = PRIV/compute/amorph/dft/make_dft.py
%   ANINIT = PRIV/compute/amorph/analysis/__init__.py        ANCERT = PRIV/compute/amorph/analysis/certify.py (docstring)
%   ANCLI  = PRIV/compute/amorph/analysis/cli.py (docstring) ANDATA = PRIV/compute/amorph/analysis/data.py (docstring)
%   ANCRY  = PRIV/compute/amorph/analysis/crystalline.py (docstring)
%   REG    = PRIV/codex/work_chat/HYPOTHESIS_REGISTER.md     LEDGER = PRIV/03_CLAIMS_LEDGER.yaml
%   DLOG   = PRIV/05_DECISION_LOG.md                         CS     = PRIV/01_CURRENT_STATE.md
%   CPLAN  = PRIV/claude/plans/COMPLETE_PLAN_2026-10.md
% ===================================================================
\subsection{First-principles calculations}
\label{sec:m-dft}

All first-principles results are \prov{COMPUTED}; the thickness laws derived from them for the device model are
\prov{FITTED} to computed points and contain no device data. % <!-- V3:58-74 -->
Every cell is neutral and spin-unpolarised, and no Hubbard correction is applied. % <!-- REG:585 ; REG:260-262 -->
Such PBE cells can show \emph{consistency} with donor behaviour; they cannot establish full ionisation or
thermodynamic charge-transition levels. % <!-- REG:260-262 -->
Hybrid-functional and spin-polarised calculations enter only as pre-declared checks.

\subsubsection{Codes, functional and pseudopotentials}
\label{sec:m-dft-codes}

All calculations use Quantum ESPRESSO \citep{Giannozzi2009,Giannozzi2017} with the PBE functional \citep{Perdew1996}
and projector-augmented-wave data sets of pslibrary 1.0.0 \citep{Bloechl1994,DalCorso2014}. The In set has the 4d
shell in the valence ($Z_{\mathrm{val}} = 13$), the W set the 5s5p shells (14); O has 6 and H 1. % <!-- RLOG:9 -->
The Al set has $Z_{\mathrm{val}} = 3$ and the Hf set 12, with 4f frozen in the core; both suggest cutoffs below the ones
used here. % <!-- SLAB:84-86 ; SLAB:145-147 -->
CPU jobs run Quantum ESPRESSO 7.5; GPU jobs run the NVIDIA build of 7.3.1 from a container
\citep{NvidiaQE,Kurtzer2017}. % <!-- DFTTEX:1352-1362 -->
On one 96-atom slab the two builds gave total energies $5\times10^{-7}$~Ry apart with the same gap and iteration count. % <!-- DFTTEX:1396-1399 -->
Structures were built with ASE \citep{Larsen2017} and verified with spglib \citep{Togo2018}; every input gives its
cell explicitly (\texttt{ibrav = 0}). % <!-- DFTTEX:600-622 -->

\subsubsection{Crystalline reference calculations}
\label{sec:m-dft-crystal}

\paragraph{Bulk.} Bixbyite \InO{} (Ia$\bar3$) \citep{Marezio1966} was converged on its 40-atom primitive cell:
71~Ry (568~Ry for the density) and a $3\times3\times3$ grid meet the pre-registered criteria of 1~mRy/atom in energy,
1~kbar in pressure and 10~meV in the gap. % <!-- DFTTEX:782 ; DFTTEX:875-879 -->
The variable-cell relaxation gives $a = \SI{10.306}{\angstrom}$, 1.87\,\% above experiment. % <!-- DFTTEX:913-915 -->
The fundamental PBE gap of the relaxed crystal, 0.8887~eV, is the reference for all confinement shifts, and the
isotropic conduction-band mass is $0.159\,m_0$. % <!-- DFTTEX:932 ; DFTTEX:942-943 -->

\paragraph{Thin films.} Hydrogen-terminated (001) slabs were built by the recipe of \citet[p.~21542]{Lin2022}:
removal of the boundary In layer, one H on each O that lost In neighbours, and \SI{25}{\angstrom} of vacuum. % <!-- DFTTEX:634-644 -->
Increasing the vacuum from 15 to \SI{35}{\angstrom} changed the electron affinity and ionisation potential by at most
0.1~meV. % <!-- DFTTEX:966-969 -->
Slabs were relaxed in atoms and in-plane cell (forces $10^{-3}$~Ry/bohr, energy $10^{-4}$~Ry) on a
$3\times3\times1$ grid with fixed occupations. % <!-- DFTTEX:857-858 -->
Band edges are read on the vacuum scale from the planar-averaged potential; the partition of the gap opening between
the conduction and valence bands uses a two-step alignment to the bulk potential, averaged over $a/2$ (the period of
bixbyite along [001]). % <!-- RLOG:213-219 -->
A projection check showed that the band edges of the relaxed \SI{2}{nm} slab are interior, not surface, states. % <!-- RLOG:192-207 -->
These calculations give the confinement shifts of claims C020 (\prov{COMPUTED}) and C021 (\prov{COMPUTED},
preliminary, termination-dependent, not yet in the results log). % <!-- LEDGER:114-124 -->

\paragraph{Tungsten in crystalline cells.} One In of the 80-atom cell was replaced by W on the 8b or the 24d site
(In$_{31}$WO$_{48}$, 3.1\,\% W), with Marzari--Vanderbilt smearing of 0.01~Ry \citep{Marzari1999}. % <!-- DFTTEX:695-697 ; DFTTEX:860 -->
W prefers the 24d site by 0.258~eV. % <!-- DFTTEX:1207-1210 -->
For the doped cells the valence-band top lies one band below that of the undoped cell, because In carries five
4d bands and W four semicore bands; an earlier analysis that missed this was corrected. % <!-- RLOG:256-263 ; CS:208 -->
The resulting W-level placement is claim C022 (\prov{COMPUTED}, crystalline cells only). % <!-- LEDGER:126-130 -->

\paragraph{Hybrid-functional check.} HSE06 \citep{Heyd2003,Krukau2006} single points at the PBE geometries used SG15
norm-conserving pseudopotentials at 80~Ry \citep{Hamann2013,Schlipf2015} and an exact-exchange cutoff of 160~Ry. % <!-- DFTTEX:1129-1136 ; DFTTEX:862 -->
The ratio of the HSE06 and PBE gap openings of the \SI{1}{nm} slab is 1.060, and 1.06--1.11 when the reduced
exact-exchange $q$ grid is taken as an uncertainty. % <!-- DFTTEX:1158-1163 -->

\paragraph{Oxygen vacancies (H-L).} Five crystalline 80-atom cells (In$_{32}$O$_{48}$, In$_{31}$WO$_{48}$,
In$_{32}$O$_{47}$ and In$_{31}$WO$_{47}$ with the vacancy near and far from W) were registered before any was
built. % <!-- REG:315 ; REG:324 ; DLOG:30 -->
The cell is fixed at $a = \SI{10.3061}{\angstrom}$; the settings are PBE, 71/568~Ry, $3\times3\times3$ for relaxation,
Gaussian smearing of 0.005~Ry, relaxation to $10^{-3}$~Ry/bohr and $10^{-4}$~Ry, W on 24d, and a
$4\times4\times4$ single point on every relaxed geometry. % <!-- REG:317-323 -->
Unfrozen numerical settings were fixed before submission (\texttt{conv\_thr} $10^{-8}$, $\beta = 0.3$, 100 SCF and
300 ionic steps). % <!-- REG:350 -->
The near site is the computed candidate with the strictly shorter W--O bond in the relaxed host (both O30 and O72 are
computed because they differ by \SI{0.0105}{\angstrom}); the far site is the O with the largest minimum-image distance
to W. % <!-- REG:345-348 ; REG:359-363 -->
Spin-polarised and 0.002~Ry single points are reported as sensitivities; a shift above 50~meV or a flipped verdict
makes the result ``not robust''. % <!-- REG:352-358 -->
The registered near cell failed one SCF; it was restarted from its best validated geometry with a pre-declared
fallback (A28). % <!-- REG:624-629 ; DLOG:40 -->
\pdPending{P16}{RESULT PENDING: H-L energies $\Delta E_{\mathrm{local}}$ and $\Delta E_{\mathrm{dop}}$, k3/k4 check,
A20 sensitivities; not yet in RESULTS\_LOG.md or the claims ledger}.

\subsubsection{Amorphous structures: machine-learned melt--quench}
\label{sec:m-dft-mace}

Amorphous structures were generated with the MACE-MPA-0 ``medium'' foundation potential
\pdPending{P17}{CITATION PENDING: MACE architecture and the MACE-MPA-0 model}, % <!-- MLV:3-5 -->
using Langevin dynamics with a \SI{1}{fs} time step and a friction of \SI{0.01}{fs^{-1}}, single precision for the
dynamics and double precision for relaxations. % <!-- MACE:124-128 ; MACE:442-443 -->
For a film, a periodic box of height $t$ at the target density is pre-relaxed, held at \SI{3000}{K} for \SI{6}{ps},
quenched linearly to \SI{300}{K} over \SI{9}{ps} and held at \SI{300}{K} for \SI{2}{ps}; a vacuum gap is then opened
along $z$, the free slab is annealed at \SI{800}{K} for \SI{3}{ps} and cooled to \SI{300}{K} over \SI{3}{ps}, relaxed
to \SI{0.05}{eV/\angstrom}, passivated and relaxed again to \SI{0.05}{eV/\angstrom}. % <!-- MACE:4-6 ; MACE:492-501 ; MACE:505-519 ; MACE:527 ; MACE:540 -->
Bulk cells skip the cleave and the passivation. % <!-- MACE:7 -->
For an interface cell the oxide is melt--quenched first (\SI{3000}{K} for \SI{4}{ps}, quench over \SI{6}{ps}) and then
held fixed while the film is melted and quenched on it. % <!-- MACE:8-10 ; MACE:473-478 -->
The seed fixes the random packing and initial velocities, but the Langevin noise comes from an unseeded global
generator, so a re-run of the same specification gives a different structure. % <!-- SLAB:293 -->
A structure enters DFT only after a quality-control gate on density, bond lengths and coordination. % <!-- CPLAN:45 ; MACE:53 -->

\subsubsection{Slab construction}
\label{sec:m-dft-slabs}

The slab rules (revision 2) were registered before any machine-learned or DFT result (decision D-14) and are applied
in code. % <!-- SLAB:1-3 ; DLOG:25 -->
\begin{itemize}
\item \emph{Doping.} W/(In+W) is exactly 2.0\,\% in every IWO cell; each W carries one extra O, which leaves one excess
electron per W (In$_{49n}$W$_n$O$_{76n}$); the undoped twin is In$_{50n}$O$_{75n}$ in the same cell. W atoms start at
least \SI{6}{\angstrom} apart. % <!-- SLAB:9-15 -->
One electron per W corresponds to about $6\times10^{20}$\cmcu{}, far above the channel carrier density of the devices. % <!-- SLAB:126 -->
\item \emph{Lateral size.} The number of W per cell is the smallest that gives an in-plane edge of at least
\SI{12.5}{\angstrom} at \SI{7.0}{g/cm^3}, so the W volume density is the same at every thickness; all cells use a
$2\times2\times1$ grid. % <!-- SLAB:22-34 -->
\item \emph{Thickness.} The label is the mass-equivalent thickness $t_{\mathrm{eq}}$ (film mass over
\SI{7.0}{g/cm^3} and area); the outer-atom span, the cation span and an electronic thickness from the planar-averaged
density are also reported. % <!-- SLAB:39-46 -->
\item \emph{Termination.} Dissociated water (OH on an under-coordinated cation, H on an under-coordinated O, in pairs,
so the electron count is unchanged), with bare-surface controls at 2.0 and \SI{0.7}{nm}. % <!-- SLAB:50-62 -->
A recombined water molecule is removed with its two H after the final relaxation and the structure is
re-relaxed. % <!-- SLAB:113-118 -->
\item \emph{Vacuum.} Every DFT slab gets exactly \SI{18}{\angstrom} of vacuum and a dipole correction along $z$. % <!-- SLAB:90-96 -->
\item \emph{Rejection.} A structure that fails quality control for any reason other than the vacuum is replaced by the
same recipe with the next seed; thresholds are never changed after a result. Three structures were replaced. % <!-- SLAB:104-111 -->
\item \emph{Doping-regime control.} A fully compensated \SI{2.0}{nm} cell with zero excess electrons (decision D-15). % <!-- SLAB:122-133 ; DLOG:24 -->
\item \emph{Interfaces.} About \SI{1}{nm} of amorphous \AlO{} (\SI{3.0}{g/cm^3}) is melt--quenched first, the film is
melt--quenched on it with the oxide fixed, the stack is cleaved below the oxide, annealed at \SI{800}{K} with all atoms
free and passivated on both outer faces. % <!-- SLAB:68-80 -->
An \HfO{} control cell uses the same recipe at \SI{9.6}{g/cm^3} (\prov{ASSUMED}) and is an alternative-interface
control only. % <!-- SLAB:139-143 ; REG:396-397 -->
Because the periodic film melt deposits film material on the outer oxide face, atoms within \SI{3}{\angstrom} of the
lowest substrate cation are excluded from every interface weight (A24, before any interface DFT result). % <!-- REG:567-574 ; SLAB:327 -->
\end{itemize}

\Cref{tab:m-dft-cells} lists the registered amorphous cells with their roles.

\begin{table}[htbp]\centering\footnotesize
\caption{Registered amorphous cells (one melt--quench sample per name). Atom counts include the passivating H where
logged. Status at the relaxation cutoff is not part of this draft.}
\label{tab:m-dft-cells}
\begin{tabularx}{\linewidth}{@{}>{\raggedright\arraybackslash}p{3.6cm}>{\raggedright\arraybackslash}p{3.2cm}>{\raggedright\arraybackslash}X@{}}
\toprule
Cell & Composition / atoms & Role and route \\
\midrule
\pdCell{aIWO_bulk_s1}, \pdCell{aIO_bulk_s1} & 50 cations, no surface & H-N bulk point; gap-anchoring references (A23); relaxed with Gaussian 0.005~Ry \\ % <!-- SLAB:35 ; REG:546-548 ; REG:408 -->
\pdCell{aIWO_2p0_s1} & In$_{98}$W$_2$ film; 270 atoms & H-G, H-N (\SI{2.0}{nm} seed 1), A15 audit, H-M reference \\ % <!-- SLAB:28 ; REG:716 ; REG:250 ; REG:409 -->
\pdCell{aIWO_2p0_s2} & In$_{98}$W$_2$O$_{156}$H$_8$; 264 & H-G, H-N (seed 2), A31 backup; Gaussian rescue (A36) \\ % <!-- REG:986 ; REG:1024 -->
\pdCell{aIO_2p0_s1}, \pdCell{aIO_2p0_s2} & In$_{100}$O$_{156}$H$_{12}$; 268 (s1) & undoped partners; A15 audit (s1) \\ % <!-- REG:581 ; REG:667 -->
\pdCell{aIWO_2p0_comp_s1} & In$_{98}$W$_2$O$_{160}$H$_{14}$; 274; $N_{\mathrm{ex}} = 0$ & H-G compensated cell; A15 audit; not an H-N cell \\ % <!-- REG:987 ; REG:1019 ; REG:1024 -->
\pdCell{aIWO_2p0_s1_bare} & bare surfaces & H-H control; parked after an SCF failure (H-H at \SI{2.0}{nm} unresolved) \\ % <!-- DLOG:46 -->
\pdCell{aIWO_1p5_s1}, \pdCell{aIO_1p5_s1} & -- & cut under the register's cut-first rule \\ % <!-- DLOG:46 ; REG:518-519 -->
\pdCell{aIWO_1p0_s1} & In$_{49}$W$_1$O$_{80}$H$_8$; 138 & H-N intermediate point; first A26 reference FINAL \\ % <!-- LEDGER:211 ; REG:604 -->
\pdCell{aIO_1p0_s1} & In$_{50}$O$_{78}$H$_6$; 134 & undoped partner; split-label qualification cell (C035) \\ % <!-- LEDGER:218 -->
\pdCell{aIWO_0p8_s2} & In$_{49}$W$_1$O$_{76}$ film; 147 & H-N; replacement of a rejected seed; Gaussian rescue (A35) \\ % <!-- REG:867 ; LEDGER:225 ; SLAB:109 -->
\pdCell{aIO_0p8_s1} & In$_{50}$O$_{75}$ film & undoped partner; A15 rescue pilot \\ % <!-- REG:412 ; SLAB:31 ; DLOG:27 -->
\pdCell{aIWO_0p7_s1}, \pdCell{aIO_0p7_s1} & In$_{49}$W$_1$O$_{83}$H$_{14}$ (W cell) & H-N thin decisive endpoint; H-H at \SI{0.7}{nm} \\ % <!-- RLOG:46 ; REG:413 ; REG:443 ; REG:266 -->
\pdCell{aIWO_0p7_s1_bare} & bare surfaces & H-H control \\ % <!-- SLAB:62 ; REG:266 -->
\pdCell{aIWO_2p0_Al2O3_s2} & 414 atoms & H-I; no feasible GPU route, stopped (H-I expected unresolved) \\ % <!-- REG:818-823 ; DLOG:46 -->
\pdCell{aIWO_1p0_Al2O3_s2} & Al$_{60}$H$_{16}$In$_{49}$O$_{174}$W$_1$; 300 & H-I; exploratory by scope (A34) \\ % <!-- REG:750 ; REG:806 -->
\pdCell{aIWO_2p0_HfO2_s1} & film of the \AlO{} cell on Hf$_{46}$O$_{92}$ & H-M (alternative-interface control) \\ % <!-- SLAB:139-141 ; REG:372-374 -->
\bottomrule
\end{tabularx}
\end{table}
\pdPending{P18}{RESULT PENDING: relaxation status of every cell at the 17 Oct 10:00 UTC cutoff (converged, exploratory,
unresolved), with segment counts, lineage and smearing family}.

\subsubsection{Validation of the machine-learned structures}
\label{sec:m-dft-mlval}

The criteria were fixed before any result: MACE forces on selected snapshots are compared with DFT single-point forces
(PBE, 60/480~Ry, Marzari--Vanderbilt 0.01~Ry, dipole correction for slabs) on the identical geometry, with limits of
\SI{0.15}{eV/\angstrom} for the force RMSE, \SI{0.10}{eV/\angstrom} for the MAE and \SI{0.25}{eV/\angstrom} for the
RMSE on W and its first-shell O; melt frames are reported, not scored. % <!-- MLV:1 ; MLV:7 ; MLV:18-26 -->
The comparator was frozen in a commit before any output was read. % <!-- LEDGER:229 ; RLOG:17 -->
The check is \textbf{failed overall} (claim C037, \prov{COMPUTED}): the undoped bulk passes (quenched RMSE 0.068, MAE
\SI{0.052}{eV/\angstrom}); the W-doped bulk fails on the W first shell only (0.434 quenched, \SI{0.449}{eV/\angstrom}
final); and the \SI{0.7}{nm} W slab fails all three limits (RMSE 0.206, MAE 0.110, W-shell
\SI{0.815}{eV/\angstrom}). % <!-- LEDGER:229 ; RLOG:20-29 -->
Most of the squared error sits in the W first shell. % <!-- LEDGER:232 ; RLOG:39 -->
Under the pre-registered consequence, the machine-learned structures are only starting points relaxed by DFT, not
validated amorphous models, and no structural statistic from the machine-learning stage alone is reported. % <!-- MLV:28-31 ; LEDGER:232 -->
That the potential's training data treat W oxides with PBE+U while the reference here is plain PBE is a possible
contributor (\prov{ASSUMED}, not tested). % <!-- LEDGER:232 ; RLOG:48 -->
\pdPending{P19}{RESULT PENDING: the two 2.0 nm force-check frames (cannot change the overall verdict; the 2.0 nm stop
rule is undetermined until they finish) and the Hf force check required by A21}. % <!-- LEDGER:229 ; REG:398-400 -->

\subsubsection{DFT relaxation}
\label{sec:m-dft-relax}

Relaxations use 60/480~Ry, Marzari--Vanderbilt smearing of 0.01~Ry for the W-doped slabs, BFGS with
$10^{-3}$~Ry/bohr and $10^{-4}$~Ry, at most 400 ionic steps, $2\times2\times1$ (slabs) or $2\times2\times2$ (bulk), a
dipole correction along $z$ for slabs, and \texttt{conv\_thr} $10^{-7}$, on GPUs with the 7.3.1 build. % <!-- MKDFT:10-11 ; MKDFT:43 ; MKDFT:54-60 ; RLOG:58 -->
The original SCF settings are $\beta = 0.2$, local-TF mixing and at most 200 iterations. % <!-- REG:799 ; REG:991 -->
Each relaxation runs as a chain of segments (an HTCondor DAG) with a wall limit of \SI{79200}{s} per segment
(\SI{165600}{s} for one cell, decision D-41). % <!-- DLOG:46 ; DLOG:52 -->
The continuation rules (decision D-24 and amendments) are:
\begin{itemize}
\item A segment restarts only from a \emph{validated} geometry, one with a converged SCF and an accepted ionic step,
and from the lowest-energy one; the restart input is a copy of the segment input in which only the atomic positions
are replaced. % <!-- REG:1139 ; REG:349 ; REG:966 -->
\item A relaxation counts as converged only if Quantum ESPRESSO prints ``bfgs converged''; it also prints final
coordinates after a failed BFGS run (A19, decision D-22). % <!-- REG:349 ; DLOG:33 -->
\item A segment ends with one of the classes converged, continue, or failed with a category (\texttt{[scf]},
\texttt{[infra]}, \texttt{[early]}, \texttt{[noprogress]}, \texttt{[parse]}); a hard failure is not replayed
(\texttt{RETRY 8 UNLESS-EXIT 2}), and a stopped segment exits with code 3 and is held for review (decision D-27). % <!-- REG:1138 ; REG:903 ; DLOG:38 -->
\item A reset after a stop follows a written procedure: the previous input is archived, the guard files are kept and
emptied, and the reset is logged in the work chat with its time, reason and input checksum. % <!-- REG:890-894 -->
\item A fatal GPU error after a new validated geometry is classified as an infrastructure failure with the geometry
archived (decisions D-37, D-38). % <!-- REG:791-797 ; DLOG:48-49 -->
\end{itemize}
To fit 250--300-atom cells on one GPU, the In and O species labels are split into two labels each with identical data
sets and settings (A25, A30, A34); at a fixed geometry the split changed the energy by $10^{-8}$~Ry and no force
component by more than $3\times10^{-8}$~Ry/bohr (claim C035, \prov{COMPUTED}). % <!-- REG:580-598 ; REG:651-658 ; LEDGER:214-219 -->
A GPU allocator setting was also permitted on this route; its retained-memory explanation is \prov{ASSUMED}. % <!-- REG:632 -->
Every final SCF restores the original labels. % <!-- REG:594 -->

\subsubsection{Smearing families and their amendments}
\label{sec:m-dft-smearing}

The W-doped slabs were registered to relax with Marzari--Vanderbilt smearing (0.01~Ry), and the undoped slabs and bulk
references relax with Gaussian smearing (0.005~Ry). The difference arose from SCF failures of undoped slabs and is
handled by amendments, each logged before the result it governs:
\begin{itemize}
\item \emph{A15 (rescue and audit).} Gaussian 0.005~Ry with $\beta = 0.1$, local-TF, 12 mixing vectors and at most
300 iterations, first as a pilot on one \SI{0.8}{nm} cell (decision D-16), then on the undoped slabs (D-26). The final
smearing is selected by an audit at fixed geometry (0.005, 0.002, 0.001~Ry on the \SI{2.0}{nm} undoped, doped and
compensated cells): edges and electron affinity within 0.05~eV, trap classes unchanged and donated charge within five
percentage points; the largest width that agrees with 0.001~Ry is selected, otherwise the quantities are unresolved.
A single-point force comparison of the two families on one \SI{2.0}{nm} doped geometry decides whether the family is a
geometry uncertainty ($10^{-3}$~Ry/bohr). % <!-- REG:1136-1146 ; DLOG:27 ; DLOG:37 -->
\item \emph{Screen at \SI{1.0}{nm}.} On the relaxed \pdCell{aIWO_1p0_s1}, switching the family changed no force
component by more than $2.4\times10^{-4}$~Ry/bohr (claim C034, \prov{COMPUTED}); this screen does not replace the
required Gaussian polish. % <!-- LEDGER:208-212 ; RLOG:169-189 -->
\item \emph{A26 (reference finals).} Gaussian 0.001~Ry reference final SCFs are permitted before the audit; the audit
rules are unchanged, self-comparison is not an audit pass, and classifications stay conditional. % <!-- REG:600-608 ; DLOG:38 -->
\item \emph{A31 (audit route and provisional width).} The nine audit runs run as CPU finals with an identical SCF path
per cell and a frozen fallback from a converged sibling width's density. Because the frozen reader does not compute
the electron affinity, the audit outcome is partial by construction and the selected width is provisional: every
dependent result carries ``A15 width provisional (EA not evaluated)''. Trap-class agreement is judged per state by the
class predicates, and named backups may substitute for a \SI{2.0}{nm} cell only on relaxation status. % <!-- REG:665-673 ; DLOG:44 -->
\item \emph{A27 controls C3 and C4.} C3 compares the 0.001~Ry reference with the audit width; C4 requires a Gaussian
0.005~Ry polish from the Marzari--Vanderbilt geometry at 2.0, 1.0 and \SI{0.7}{nm}, and a crossover must hold on the
polished geometries. % <!-- REG:493-501 ; REG:459-466 -->
\item \emph{A35 (\pdCell{aIWO_0p8_s2}).} After three SCF failures this cell continues from its best validated geometry
with Gaussian 0.005~Ry and the rescue aids, one attempt only; its class is read on a Gaussian-family geometry. % <!-- REG:866-903 ; REG:933-934 ; DLOG:47 -->
\item \emph{A36 (\SI{2.0}{nm} doped chains).} \pdCell{aIWO_2p0_s2} continues with Gaussian 0.005~Ry plus the aids and
\pdCell{aIWO_2p0_comp_s1} with Marzari--Vanderbilt 0.01~Ry plus the aids, one attempt each; \pdCell{aIWO_2p0_s1} gets
the Gaussian route once only if it writes an SCF failure. % <!-- REG:984-1015 ; DLOG:51 -->
\end{itemize}
The W series therefore uses mixed relaxation smearing, and every table of H-G, H-M and H-N gives each cell's family
and lineage. % <!-- REG:939-943 ; REG:1025 -->
Because a Gaussian-rescued \pdCell{aIWO_2p0_s2} has no Marzari--Vanderbilt geometry, control C4 cannot be satisfied for
it, and the H-N \SI{2.0}{nm} endpoint is \prov{UNRESOLVED}; neither a crossover nor its absence can be claimed from an
A36 geometry. % <!-- REG:1016-1018 -->
\pdPending{P20}{RESULT PENDING: A15 audit outcome and the provisional width (A31); whether aIWO\_2p0\_s1 needed the
conditional A36 rescue; C4 polishes at 1.0 and 0.7 nm}.

\subsubsection{Restart reproducibility and qualification of a relaxation}
\label{sec:m-dft-restart}

For \pdCell{aIWO_0p8_s2} a fresh SCF at a fixed geometry with identical settings did not reproduce the continued SCF
at production tolerance: the energy differed by $1.18\times10^{-5}$~Ry and a force component by
$4.99\times10^{-3}$~Ry/bohr (claim C036, \prov{COMPUTED}). % <!-- LEDGER:221-226 ; RLOG:77-84 -->
Whether this is an incomplete-SCF force error or two converged electronic branches is open; the branch and dipole
mechanism is \prov{ASSUMED}. Numerical-path invariance is therefore not assumed for this cell, and nothing is claimed
for other cells. % <!-- REG:970-976 ; LEDGER:225 -->
For every cell on a rescue route (A35, A36), two gates follow:
\begin{itemize}
\item \emph{Restart gate.} At every later segment restart the first, fresh SCF must reproduce the previous segment's
value at the same geometry within $10^{-5}$~Ry in energy and $10^{-3}$~Ry/bohr in every force component; a failed gate
does not stop the chain. % <!-- REG:911-916 ; REG:1021 -->
\item \emph{Qualification SCF.} After convergence, a fresh SCF at the final geometry with the relaxation settings must
reproduce the relaxation's last SCF within the same tolerances. % <!-- REG:917-921 ; REG:1022 -->
\end{itemize}
The gate evidence is bound with sha256 into the cell's run certificate. A failed or missing qualification, or any
failed gate on the lineage, makes the relaxation \emph{exploratory}; the frozen reader then excludes the cell from
primary scoring, and the flag is inherited by every geometry relaxed from it. % <!-- REG:922-927 ; REG:1023 -->
Independently, any relaxation not converged at the 17 October 10:00~UTC cutoff is exploratory, and a primary test
without a converged geometry is unresolved. % <!-- REG:1076 -->
\pdPending{P21}{RESULT PENDING: restart-gate and qualification-SCF outcomes for aIWO\_0p8\_s2, aIWO\_2p0\_s2 and
aIWO\_2p0\_comp\_s1}.

\subsubsection{Final SCF, export and certification}
\label{sec:m-dft-final}

Final SCFs run on CPUs with Quantum ESPRESSO 7.5 at 71/568~Ry (the cutoff of the crystalline study), the
$2\times2\times1$ grid and $1.3\times$ the occupied band count, followed by the planar-averaged potential and
projections (\texttt{projwfc.x}); they are submitted on hold until their smearing gate is cleared. % <!-- MKDFT:12-14 ; DLOG:29 -->
The reference finals use Gaussian 0.001~Ry with one $k$-point pool and the registered SCF settings
(\texttt{conv\_thr} $10^{-7}$, $\beta = 0.2$, 200 iterations, local-TF). % <!-- REG:604-606 ; DLOG:50 -->
Before scratch space is lost, the following are exported: per-state projections with atom and orbital maps and
completeness values, projected densities of states, occupations and $k$ weights, planar averages, final geometries
with their lineage, eigenvalues, the actual inputs and logs of every segment, and the charge density of the
\SI{2.0}{nm} cells (A14). % <!-- REG:1078-1086 -->
A corrected export (decision D-24, item C) was required before any final was released, and completeness is verified
against this list, not inferred from job success. % <!-- DLOG:35 ; REG:607-608 -->

The read-out is a frozen Python package (version 1.0.0-rc7; code sha256 \texttt{fff642a1}\ldots\texttt{e4f01158};
221 tests), frozen after Codex reported it ready and before any slab read-out existed (decision D-32). % <!-- DLOG:43 ; ANINIT:8 -->
It reads the registered export (the XML eigenvalue and projection files); a text route with rounded projections is
used only for the crystalline reproduction. % <!-- ANDATA:4-10 -->
Before use on amorphous cells it must reproduce two crystalline numbers: a 3\,\% W-5d share and
$\dEc = \SI{0.281}{eV}$ at \SI{2}{nm}. % <!-- REG:1068 ; ANCRY:1-2 -->
\pdPending{P22}{CHECK PENDING: record of the crystalline reproduction by the frozen reader (register section 3)}.
For each final the DFT agent compiles a certificate holding the sha256 of every file the read-out consumes, the
geometry and pseudopotential hashes, a composition manifest and one block per qualification (relaxation, smearing,
geometry polish) with sha256-bound evidence; an automated gate recomputes every hash, and Codex audits the
certificate. Without a valid certificate every primary score is unresolved. % <!-- ANCERT:4-17 ; ANCLI:3-6 -->
Every output carries the code, register, input and certificate hashes and a self-hash. % <!-- ANCERT:17 ; ANCLI:19-21 -->

\subsubsection{Read-out definitions}
\label{sec:m-dft-readout}

The definitions were frozen with the register (section 1, amended by A16 and A24 before any electronic result). % <!-- REG:43-93 -->
\begin{itemize}
\item \emph{Projections.} The weight of atom $i$ in state $(n,k)$ is its L\"owdin projection \citep{Lowdin1950} summed
over orbitals and normalised over atoms; a state with completeness below 0.90 is flagged, and a test resting on flagged
states is unresolved. % <!-- REG:45-49 -->
\item \emph{Localisation.} $\mathrm{PR}_n = 1/(N\sum_i p_i^2)$ over all $N$ atoms of the cell; extended if
$\mathrm{PR}_n \ge 0.20$, localised if $\le 0.05$. % <!-- REG:50-55 -->
\item \emph{Occupations.} Each state contributes $f_{nk}w_k$, with $f_{nk}$ from the final SCF's own smearing
function. % <!-- REG:56-59 -->
\item \emph{Host edge.} The lowest eigenvalue above band $N_v$ of any extended, interior-weighted state (interior
weight at least 0.8 times the interior atom fraction); unresolved if none qualifies within 1.0~eV or the interior is
empty. A host valence edge is defined in the same way at or below $N_v$. % <!-- REG:60-67 ; REG:80 -->
\item \emph{Electron accounting.} $N_v = (N_{\mathrm{el}} - N_{\mathrm{ex}})/2$ with $N_{\mathrm{ex}} = n_W$ in the
W-doped cells and 0 in the compensated and undoped cells; the occupation above $N_v$ must equal $N_{\mathrm{ex}}$
within 0.02~e, and every fraction uses the nominal denominator. % <!-- REG:72-82 -->
\item \emph{Alignment.} Valid only with an interior plateau of the planar-averaged potential with a standard deviation
of at most 0.05~eV. % <!-- REG:83-85 -->
\end{itemize}
H-N classifies each W-doped cell from these quantities as DONOR ($f_{\mathrm{ext}} \ge 0.80$, $\Delta_W > 0.3$~eV,
$f_{\mathrm{locW}} < 0.10$), ECB ($f_{\mathrm{ext}} < 0.50$ or $f_{\mathrm{locW}} \ge 0.50$), PARTIAL or UNRESOLVED;
these are \prov{ASSUMED} operational screens, not universal donor criteria. % <!-- REG:427-432 -->
\pdPending{P23}{RESULT PENDING: H-N per-cell classes and the series outcome; H-G, H-H, H-K scores; release decision
D-42 (not yet in the decision log)}.

\subsubsection{Gap anchoring and hybrid-functional checks}
\label{sec:m-dft-gap}

Absolute gaps are not taken from PBE. The scheme $E_g(t) = E_{g,\mathrm{anchor}} + r\,\Delta E_g^{\mathrm{PBE}}(t)$
uses a measured anchor (3.05~eV, \prov{LITERATURE}, a different sputtered IWO device \citep{Fan2021}), the host-edge
gap of each slab minus that of the matching amorphous bulk, and the HSE06/PBE ratio $r$ = 1.06--1.11 with its spread as
an uncertainty (A23). % <!-- REG:540-544 -->
Two hybrid checks were registered with read-out rules: a crystalline W cell (is the W-5d state still above the host
edge in HSE06?) and an amorphous bulk (does its HSE06--PBE gap offset differ from the crystalline 1.188~eV by more than
0.10~eV?). % <!-- REG:550-563 -->
\pdPending{P24}{RESULT PENDING: A23 (ii) W-cell and (iii) amorphous-bulk hybrid checks; amorphous bulk references}.

\subsubsection{Computing resources}
\label{sec:m-dft-resources}

The amorphous and vacancy calculations ran on CERN's HTCondor batch system, submitted from lxplus, with inputs and
results on EOS \citep{Thain2005}. % <!-- DFTTEX:1352-1357 ; CPLAN:94 -->
Relaxations used single H100 NVL GPUs (A100 allowed for the vacancy cells); the H100/H200 pool was saturated, two-GPU
requests matched no slot, and H200 nodes were not available to the group. % <!-- DLOG:34 ; REG:678-680 -->
The \SI{2.0}{nm} chains used 85.6--90.2~GB of the 95.8~GB of GPU memory on the one-GPU route. % <!-- REG:715 ; REG:1007 -->
Kerberos tickets of running jobs are not renewed, so long jobs need a separate credential-refresh job. % <!-- CS:307 ; DLOG:46 -->
The registered deadlines were: relaxations by 17~October 10:00~UTC, single-point controls by 20~October 12:00~UTC, last
submission 18~October 22:00~UTC, access ending 21~October; the observed throughput was about 0.8 GPU starts per hour,
an observation, not a guaranteed capacity. % <!-- REG:530-532 ; REG:1067-1087 -->
The crystalline stage used about 8300 weighted CPU core-hours and 85--90 GPU-hours. % <!-- DFTTEX:1517-1518 -->
\pdPending{P25}{RESULT PENDING: total CPU and GPU hours of the amorphous, vacancy and validation campaign}.

% ===================================================================
% methods_registration.tex -- Methods: governance, the hypothesis register (rules, freeze tags), the amendment log
% A1-A36 in summary, sight history and disclosures.
% Draft 2026-10-13 (Claude), for human review; no claim here is approved for a manuscript (charter section 9).
% SIGHT RULE: sight events are described by what was seen and when, never by the values seen.
% ASCII only. Every number carries a source comment  % <!-- KEY:line -->  (file:line). KEY legend:
%   PRIV    = iwo-tft-research-private/                     PKG    = IWO_PHYSICS_CONSTRAINED_MODEL_V1/
%   CHARTER = PRIV/00_PROJECT_CHARTER.md                    CS     = PRIV/01_CURRENT_STATE.md
%   LEDGER  = PRIV/03_CLAIMS_LEDGER.yaml                    DLOG   = PRIV/05_DECISION_LOG.md
%   REG     = PRIV/codex/work_chat/HYPOTHESIS_REGISTER.md   SLAB   = PRIV/compute/amorph/SLAB_PROTOCOL.md
%   MLV     = PRIV/claude/plans/ML_VALIDATION.md            NPLAN  = PRIV/claude/NATURE_PLAN.md
%   CD      = PKG/analysis_2026-10-10_campaign/CAMPAIGN_DESIGN.md
%   S2PREP  = PKG/analysis_2026-10-10_campaign/S2_PREP.md   S2FIX  = PKG/analysis_2026-10-10_campaign/S2_SCORER_FIXES_2026-10-12.md
%   S2GATE  = PKG/analysis_2026-10-10_campaign/S2_RUNNER_GATES_2026-10-13.md
%   RLOG    = PKG/qe_workflow/cern_htcondor/results/RESULTS_LOG.md   V2TEX = PKG/report_latex_full/chapters/14_model_v2.tex
%   GIT     = local git tag records of PRIV (read 2026-10-13 with git for-each-ref; not a file:line source)
% ===================================================================
\subsection{Pre-registration, amendments and sight history}
\label{sec:m-reg}

\subsubsection{Governance}
\label{sec:m-reg-governance}

The work follows a written charter. Every scientific statement carries one status label (for example
\prov{MEASURED}, \prov{DERIVED}, \prov{FITTED}, \prov{COMPUTED}, \prov{HYPOTHESIS}, \prov{UNRESOLVED},
\prov{FALSIFIED}), and labels change only with new evidence, recorded in a claims ledger. % <!-- CHARTER:27-39 -->
No target-device-specific post-hoc fitting is allowed: global parameters may be calibrated on a declared calibration
set, and anything tuned after a holdout has been revealed is calibration, never prediction. % <!-- CHARTER:41-47 -->
A prediction is registered against a frozen model with numerical pass/fail criteria and timestamped outside the
repository before any holdout is released. % <!-- CHARTER:86-95 -->
The AI agents had defined roles: Claude Code as computational lead and record keeper, Codex as independent reviewer
and hypothesis challenger; agreement between AI systems is not counted as evidence, and the humans decide. % <!-- CHARTER:15-21 ; CHARTER:53 ; CHARTER:105-116 -->
\pdPending{P26}{DECISION PENDING: the charter is still marked PROPOSED and its adoption (decision D-10) is pending; adopt
it or describe the governance as practised}. % <!-- CHARTER:7 ; DLOG:21 -->

No device data were withheld in this project: every measured curve was seen before any model of this paper was built.
The registrations therefore protect the analysis and model choices against the outcome of our own calculations; they
are not blind tests against unseen measurements. % <!-- REG:37 ; CD:18 ; CD:501 ; DLOG:22 -->

\subsubsection{Registration record}
\label{sec:m-reg-record}

\Cref{tab:m-reg-record} lists the registered documents in time order.

\begin{table}[htbp]\centering\footnotesize
\caption{Registered documents and freezes. Times are UTC unless stated; decision-log dates are local dates (UTC+8).}
\label{tab:m-reg-record}
\begin{tabularx}{\linewidth}{@{}>{\raggedright\arraybackslash}p{3.0cm}>{\raggedright\arraybackslash}X>{\raggedright\arraybackslash}p{3.6cm}@{}}
\toprule
When & What was frozen & Record \\
\midrule
before the V2 launches & V2 pass/fail criteria and predictions; one correction before the held-out outcome & commits \texttt{7b75ee6}, \texttt{a914165}; outcome \texttt{852dc8a} \\ % <!-- CS:164 ; V2TEX:43-50 -->
before any comparison & extraction protocol; amendment A1 (normalised SS window) & commits \texttt{7a61687}, \texttt{7ce7116} (D-05) \\ % <!-- CS:86-87 ; DLOG:16 -->
2026-10-10, before any result & slab protocol revision 2 (D-14); force-validation criteria for the machine-learned structures & \file{SLAB_PROTOCOL.md}; \file{ML_VALIDATION.md} \\ % <!-- SLAB:1 ; DLOG:25 ; MLV:1 -->
2026-10-10, before any launch & ATLAS campaign design (C10) & \file{CAMPAIGN_DESIGN.md} \\ % <!-- CD:1 ; CD:12 -->
2026-10-10T18:33:13Z & hypothesis register, revision 2 (D-17) & tag \texttt{register-rev2-frozen-}\allowbreak\texttt{20261010T1833Z} \\ % <!-- REG:3 ; DLOG:28 ; GIT: creator date 2026-10-11 02:33:15 +0800 -->
2026-10-11T09:09Z, before any slab electronic result was read & H-N read-out rules and controls (A27, D-28) & register A27 \\ % <!-- REG:610 ; DLOG:39 -->
2026-10-11T22:49Z, before any slab read-out & DFT read-out code 1.0.0-rc7, 221 tests (D-32) & tag \texttt{readout-1.0.0-}\allowbreak\texttt{rc7-frozen-}\allowbreak\texttt{20261011T2249Z} \\ % <!-- DLOG:43 ; GIT -->
before any force-check output was read & force comparator v1.0.0 & commit \texttt{d0179e1} \\ % <!-- LEDGER:229 ; RLOG:17 -->
2026-10-12, before S2-11 was unsealed & TCAD register scorer (sha256 \texttt{c0f28bb9}\ldots) & \file{S2_SCORER_FIXES_}\allowbreak\file{2026-10-12.md} \\ % <!-- S2FIX:1 ; S2FIX:349 ; S2GATE:243 -->
\bottomrule
\end{tabularx}
\end{table}
\pdPending{P27}{CHECK PENDING: external timestamp of the register and read-out freezes (charter step 5: a Zenodo/OSF DOI,
or at least the tags pushed to GitHub); the remote of the record is private}. % <!-- CHARTER:92 -->

\subsubsection{The hypothesis register}
\label{sec:m-reg-register}

The register distinguishes a \emph{registered prediction} (a frozen pass/fail test) from a \emph{forecast} (a
drafter's estimate, recorded but not tested); every threshold is an \prov{ASSUMED} decision threshold. % <!-- REG:13-15 -->
Under its censoring rule a quantity below the demonstrated numerical resolution is \emph{below resolution}, counts as
consistent with an upper bound only if that resolution lies below the bound, and is otherwise \emph{unresolved}, never
a pass. % <!-- REG:16-20 -->
Section~1 of the register freezes the read-out definitions (\cref{sec:m-dft-readout}); section~3 the experiment plan,
deadlines and export list; section~4 lists what is frozen, including the scoring windows, stop rules, calibration
membership, seed exclusions and the amendment log. % <!-- REG:43-93 ; REG:1030-1095 -->
After the freeze every change is a logged amendment, and the frozen predictions are scored and reported whatever the
outcome. % <!-- REG:3 ; DLOG:28 -->
\Cref{tab:m-reg-hyp} lists the hypotheses.

\begin{table}[htbp]\centering\footnotesize
\caption{Hypotheses of the register. ``Registered'': frozen pass/fail predictions; ``exploratory'': frozen read-out
rules, no direction scored.}
\label{tab:m-reg-hyp}
\begin{tabularx}{\linewidth}{@{}l>{\raggedright\arraybackslash}X>{\raggedright\arraybackslash}p{4.0cm}@{}}
\toprule
ID & Hypothesis (short) & Type; method \\
\midrule
H-A & The \SI{2.0}{nm} floor is extrinsic or instrument-limited & registered; decisive part expected unresolved (no instrument records) \\ % <!-- REG:99 ; REG:114-119 -->
H-B & One shared ungated strip explains the floor scaling & registered; TCAD \\ % <!-- REG:100 ; REG:121-135 -->
H-C & Gate leakage and generation do not set the floor & registered; TCAD \\ % <!-- REG:101 ; REG:137-158 -->
H-D & Band mobility needs $1/t$ scattering, not only the mass & registered; TCAD (mass-only part not computed, A22) \\ % <!-- REG:102 ; REG:160-177 -->
H-E & Confinement explains part of $\Vth(t)$ & registered (amended by A13); TCAD \\ % <!-- REG:103 ; REG:179-236 -->
H-F & The \SI{6.3}{nm} shape residual is unexplained & registered; TCAD (scoring definitions pending) \\ % <!-- REG:104 ; REG:238-247 ; S2PREP:167-173 -->
H-G & W behaves as a donor into host states, no pinning state & registered; DFT, \SI{2.0}{nm} \\ % <!-- REG:105 ; REG:249-263 -->
H-H & OH termination moves dipoles, not the interior edge & registered; DFT \\ % <!-- REG:106 ; REG:265-271 -->
H-I & \AlO{} changes the alignment without channel-edge traps & registered; DFT \\ % <!-- REG:107 ; REG:273-284 -->
H-J & One shared model transfers across 2.0, 6.3, \SI{13.2}{nm} & registered; TCAD \\ % <!-- REG:108 ; REG:286-302 -->
H-K & Amorphous confinement agrees with the crystalline laws & registered; DFT \\ % <!-- REG:109 ; REG:304-312 -->
H-L & W suppresses O vacancies in crystalline \InO{} & registered after the freeze (A17), before any cell was built; DFT \\ % <!-- REG:110 ; REG:315 -->
H-M & \HfO{} interface control at \SI{2.0}{nm} & registered after the freeze (A18), before any structure; DFT \\ % <!-- REG:111 ; REG:370 -->
H-N & Extended-carrier efficacy of W against thickness & exploratory after the freeze (A27), before any slab electronic result; DFT \\ % <!-- REG:112 ; REG:405 ; DLOG:39 -->
\bottomrule
\end{tabularx}
\end{table}

\subsubsection{Amendment log}
\label{sec:m-reg-amendments}

Amendments A1--A16 turned revision~1 into the frozen revision~2, in response to two Codex reviews, a numerical audit of
the effective density of states and Codex's confirmation. % <!-- REG:5-9 ; REG:1097-1116 -->
A17--A36 were added after the freeze; A32 and A33 are reserved for two pilots that have not been registered. % <!-- REG:4 ; REG:676 -->
Each post-freeze amendment was approved by the user and recorded in the decision log, and most after a Codex review
with requested changes applied. \Cref{tab:m-reg-amend} gives the post-freeze amendments with their timing relative to
the results they govern.

\begin{table}[htbp]\centering\footnotesize
\caption{Amendments after the freeze (register section 2 and section 5). Times UTC.}
\label{tab:m-reg-amend}
\begin{tabularx}{\linewidth}{@{}l>{\raggedright\arraybackslash}p{2.6cm}>{\raggedright\arraybackslash}X@{}}
\toprule
\# & Logged & Content (timing) \\
\midrule
A17 & 10-10 18:48 & H-L added, before any vacancy cell was built or submitted (D-19) \\ % <!-- REG:315 ; REG:1117 ; DLOG:30 -->
A18 & 10-10 18:52 & H-M added, before the Hf data set was installed or any structure built (D-20) \\ % <!-- REG:370 ; REG:1118 ; DLOG:31 -->
A19 & 10-10 19:26 & H-L near/far site rule; ``bfgs converged'' convergence rule; settings fixed before submission \\ % <!-- REG:344 ; REG:1119 -->
A20 & 10-10 20:42 & H-L robustness sensitivities and ``not robust'' rule, before any H-L result \\ % <!-- REG:352 ; REG:1120 -->
A21 & 10-10 20:45 & H-M limited to an alternative-interface control; Hf force check; timing erratum (a non-production smoke build existed) \\ % <!-- REG:396-402 ; REG:1121 -->
A22 & 10-10 20:45 & amorphous dispersion workflow descoped; mass-only H-D prediction reported ``unresolved: not computed'' \\ % <!-- REG:170-177 ; REG:1122 -->
A23 & 10-10 21:06 & gap anchoring; amorphous bulk references; two hybrid checks with read-out rules (D-25) \\ % <!-- REG:539 ; REG:1123 ; DLOG:36 -->
A24 & 10-11 03:13 & interface recipe kept; outer oxide face excluded from interface weights; crystallinity flags, before any interface DFT result (D-26) \\ % <!-- REG:567 ; REG:1124 ; DLOG:37 -->
A25 & 10-11 08:14 & split species labels for one \SI{2.0}{nm} cell on one GPU, with energy gate and force qualification (D-27) \\ % <!-- REG:580 ; REG:1125 -->
A26 & 10-11 08:14 & Gaussian 0.001~Ry reference finals before the smearing audit; classifications conditional (D-27) \\ % <!-- REG:600 ; REG:1126 -->
A27 & 10-11 09:09 & H-N with classes, outcome rules and controls C1--C6, before any slab electronic result (D-28) \\ % <!-- REG:610 ; REG:1127 -->
A28 & 10-11 10:06 & H-L near-cell restart and pre-declared fallback; erratum on the number of validated SCFs (D-29) \\ % <!-- REG:624-629 ; REG:1128 -->
A29 & 10-11 12:53 & GPU allocator setting for one cell (cause \prov{ASSUMED}) (D-30) \\ % <!-- REG:631 ; REG:1129 -->
A30 & 10-11 20:30 & one-GPU split-label route for seven idle chains; erratum on the out-of-memory path (D-31) \\ % <!-- REG:651 ; REG:659-663 ; REG:1130 -->
A31 & 10-12 00:16 & smearing-audit route, SCF fallback, provisional width without the electron affinity, role substitution, before any \SI{2.0}{nm} final (D-33) \\ % <!-- REG:665 ; REG:1131 -->
A34 & 10-12 06:06 & one-GPU routes for the \SI{2.0}{nm} force-check frames and one interface chain; recovery-rule correction; the interface route exploratory by scope (D-34, D-37) \\ % <!-- REG:675 ; REG:1132 ; DLOG:45 ; DLOG:48 -->
A35 & 10-12 07:27 & Gaussian rescue of \pdCell{aIWO_0p8_s2}; restart gate and qualification SCF; exploratory rule; post-hoc disclosure (D-36) \\ % <!-- REG:866 ; REG:1133 ; DLOG:47 -->
A36 & 10-12 13:21 & rescues of two \SI{2.0}{nm} W chains and a conditional third; H-N \SI{2.0}{nm} endpoint unresolved without a C4 polish; sight erratum (D-40, D-41) \\ % <!-- REG:984 ; REG:1134 ; DLOG:51-52 -->
\bottomrule
\end{tabularx}
\end{table}

Two pre-freeze amendments matter for the reading of the TCAD results. A13 added a \SI{2.0}{nm} confinement band that
was informed by the S1 results, before any S3 run; the original band is preserved and both are scored, and the failed
S1-05 expectation stays recorded as a failure. A15 introduced the smearing rescue and audit
(\cref{sec:m-dft-smearing}). % <!-- REG:1113 ; REG:1115 ; REG:184-190 ; REG:211-212 -->

\subsubsection{Sight history and disclosures}
\label{sec:m-reg-sight}

Sight events are listed by what was seen and when; no value seen is reported here.
\begin{enumerate}
\item \emph{Before the freeze.} The smoke runs S0-01 and S0-02 and all of stage S1 had been seen. At the freeze the
team was blind to S2--S7 and to every amorphous electronic target quantity; it had seen the SCF-convergence failures of
two undoped slabs (convergence behaviour, one first-SCF energy and one force norm). % <!-- REG:21-36 -->
\item \emph{Before A27.} No slab final had run. The electronic outputs seen were relaxation-SCF energies and iteration
counts, and geometry-selection summaries that included Kohn--Sham gaps of \pdCell{aIO_0p8_s1}. % <!-- REG:615 -->
\item \emph{After A27.} The results log of the split-label qualification records a Fermi energy of
\pdCell{aIO_1p0_s1} at a non-final geometry. \pdPending{P28}{DISCLOSURE PENDING: this entry is not listed in the
register's sight history; decide whether it needs a register erratum}. % <!-- RLOG:118-145 ; LEDGER:214-219 -->
\item \emph{A35.} The relaxation outputs of segments s0--s2 of \pdCell{aIWO_0p8_s2} were read for energies, energy
components, forces and SCF accuracy. Diagnostic scripts also parsed relaxation-SCF Fermi energies, which appeared in
a review digest but were used for no conclusion or decision. % <!-- REG:944-948 ; RLOG:113-115 ; LEDGER:225 -->
\item \emph{A36 erratum: partial electronic unblinding.} On 2026-10-12, before A36 was committed, a research agent
parsed relaxation-SCF Fermi energies and Kohn--Sham band eigenvalues of four \SI{2.0}{nm} cells
(\pdCell{aIO_2p0_s1}, \pdCell{aIWO_2p0_s2}, \pdCell{aIWO_2p0_comp_s1} and a thread-test SCF of
\pdCell{aIWO_2p0_s1}) to derive the position of one high band relative to the Fermi level for a band-count proposal
(team exposure about 12:49Z); the main session read the four derived values at about 13:28Z. No projections,
participation ratios, gaps or edges were computed; the values informed no decision, the band-count change was
dropped, and the files are quarantined with sha256. Every \SI{2.0}{nm} H-G and H-N table therefore carries a linked
disclosure footnote; no statement that no electronic quantity was seen is made; the unchanged projection,
participation-ratio and class tests are not automatically invalidated; and later control cuts or scoring amendments
acknowledge this history and preserve the original unresolved scores. % <!-- REG:1027 ; DLOG:52 -->
\item \emph{Procedural disclosure (A35).} After a pre-declared condition for switching \pdCell{aIWO_0p8_s2} to Gaussian
smearing had been met, the chain was released once more on the original smearing on a user go; this was recorded only
in a CERN log, not in the work chat as the reset procedure requires, and was disclosed after the fact. % <!-- REG:950-959 ; DLOG:47 -->
\item \emph{Restart non-reproducibility.} The finding of claim C036 led to the restart gate and qualification SCF of
\cref{sec:m-dft-restart}. % <!-- LEDGER:221-226 ; REG:960-976 -->
\item \emph{TCAD.} The deviations of \cref{sec:m-tcad-controls}: the S0-02 repair cap, launches while the regression
gate was red (R23), and the sealed S2-11 run with its pending ``possibly exposed'' ruling. % <!-- REG:39 ; S2PREP:82 ; S2GATE:123 -->
\item \emph{Errata.} A21 (timing of a smoke build), A28 (number of validated SCFs) and A30 (out-of-memory path) were
corrected in the register without changing any rule. % <!-- REG:401-403 ; REG:629 ; REG:659-663 -->
\end{enumerate}
\pdPending{P29}{TEXT PENDING: final sight-history table per cell (what was read, by whom, when, before or after each
freeze) once all finals are released; release decision D-42}.

\subsubsection{Reporting commitments}
\label{sec:m-reg-reporting}

Every frozen prediction is reported with its score, including failures and unresolved outcomes. Classes read on an
exploratory geometry are reported but never count as resolved coverage; registered unresolved scores are not
replaced by later analyses. The \SI{13.2}{nm} swing mismatch and every failed branch and launch are reported (decision
D-06). Calibration is never called validation. % <!-- DLOG:28 ; REG:927 ; REG:1027 ; DLOG:17 ; CD:18 ; CD:511 -->
\pdPending{P30}{RESULT PENDING: the scorecard of every registered prediction (H-A to H-M) and the H-N exploratory
read-out, each with its unresolved checks}.


% =====================================================================================================================
\section{Limitations}
% ===================================================================
% limitations.tex -- Limitations of the computational study (data, device model, first-principles calculations,
% procedure). Draft 2026-10-13 (Claude), for human review; no claim here is approved for a manuscript (charter section 9).
% ASCII only. Every number carries a source comment  % <!-- KEY:line -->  (file:line). KEY legend:
%   PRIV    = iwo-tft-research-private/                     PKG    = IWO_PHYSICS_CONSTRAINED_MODEL_V1/
%   CHARTER = PRIV/00_PROJECT_CHARTER.md                    CS     = PRIV/01_CURRENT_STATE.md
%   LEDGER  = PRIV/03_CLAIMS_LEDGER.yaml                    DLOG   = PRIV/05_DECISION_LOG.md
%   REG     = PRIV/codex/work_chat/HYPOTHESIS_REGISTER.md   SLAB   = PRIV/compute/amorph/SLAB_PROTOCOL.md
%   MLV     = PRIV/claude/plans/ML_VALIDATION.md            V3     = PKG/config/iwo_material_model_v3.yaml
%   CD      = PKG/analysis_2026-10-10_campaign/CAMPAIGN_DESIGN.md
%   S2PREP  = PKG/analysis_2026-10-10_campaign/S2_PREP.md   S2FIX  = PKG/analysis_2026-10-10_campaign/S2_SCORER_FIXES_2026-10-12.md
%   V2TEX   = PKG/report_latex_full/chapters/14_model_v2.tex
% ===================================================================
\label{sec:limitations}

\paragraph{Measured data.}
\begin{itemize}
\item One device per thickness, one drain bias (\SI{0.7}{V}), no gate current and no reverse sweep: device-to-device
spread cannot be quoted, and no thickness trend can be separated from it. % <!-- LEDGER:5-9 ; LEDGER:26 -->
\item The thicknesses are nominal, and whether the devices come from one fabrication run, and their sweep history,
are unknown. % <!-- CS:253-254 -->
\item The \SI{2.0}{nm} off-state floor (\SI{1.3}{pA} in total) is at the level of the measurement set-up. % <!-- LEDGER:20 -->
\item The \SI{31.8}{nm} device never turns off; it is a documented model failure and not a holdout. % <!-- LEDGER:172-176 ; CS:245-246 -->
\item The source paper's mobility and near-thermal swing differ from ours on the same curves; a definition difference
reconciles the mobility numerically but the paper's recipe is unconfirmed (claims C024, C025 \prov{UNRESOLVED}). % <!-- LEDGER:138-152 -->
\item No new measurements were possible (decision D-11). The measurement that separates an ohmic parallel strip
(floor $\times7$ from 0.1 to \SI{0.7}{V}) from a diffusion form ($\times1.02$) could not be made, so the TCAD floor
tests bound and rank mechanisms only within the model. % <!-- DLOG:22 ; CS:219-220 ; CD:395 -->
\item Which oxide touches the channel is not confirmed by the fabrication group; the model places \AlO{} against the
channel. % <!-- CS:241 ; SLAB:68-70 -->
\item Every comparison with device data is retrospective; leave-one-film-out tests measure transferability within four
already-seen devices, not independent prediction. % <!-- REG:37 ; CD:501 -->
\end{itemize}

\paragraph{Device model.}
\begin{itemize}
\item The parameters are not unique: work function, affinity, fixed charge and conduction-band shift act as one number,
the band mobility trades off against the tail density, and series resistance against the mobility exponent (claim
C018). Confinement is physically expected but cannot be identified from transfer curves. % <!-- LEDGER:104-107 ; LEDGER:118 ; CS:153-155 -->
\item Every out-of-sample test of the \SI{6.3}{nm} threshold has failed (claim C014, \prov{FALSIFIED}); the earlier
models are calibrated descriptions, not predictive models of the thickness dependence. % <!-- LEDGER:80-84 ; V2TEX:334-335 -->
\item The earlier calibrated model did not reproduce the gate-voltage dependence of $\muFE$ at 2.0 and \SI{6.3}{nm}, and
its \SI{13.2}{nm} swing is too large (claims C012, C010). % <!-- LEDGER:57-61 ; LEDGER:69-72 -->
\item The quasi-two-dimensional density of states at \SI{2.0}{nm} is a substitution inside the 3-D Fermi--Dirac form; it
overstates the degenerate on-state density by up to $1.7\times$, and the on-current at \SI{3}{V} by an inferred
7--8\,\%. The density-of-states convention cannot be told apart from fixed charge, tail density and Gaussian depth in
subthreshold data. % <!-- CD:111 ; REG:226-228 -->
\item The rigid-shift interpretation of the confinement tests holds only because trap levels follow $\Ec$ in the model
(\prov{ASSUMED}). % <!-- REG:206 -->
\item The DFT laws are extrapolated beyond \SI{2}{nm}, the HSE06/PBE ratio was measured at \SI{1}{nm} only, and the
device thickness is identified with the slab H--H thickness (\prov{ASSUMED}). % <!-- CD:151 ; V3:76 -->
\item Contacts are ideal Ohmic and unconstrained (no transfer-length data); the solution is electron-only; the donor
density is constant by assumption. % <!-- V3:209-211 ; V3:124-125 -->
\item The leakage-maximising oxide set is an optimistic scenario, not a bound: pinholes, local thinning and percolation
paths are not modelled. The oxide trap-assisted tunnelling runs are not convergence-qualified. % <!-- CD:376 ; REG:157 ; S2PREP:164 -->
\item The numerical-resolution method of the scorer is an \prov{ASSUMED} operational method; with the current run plan
most primary leakage and affinity-only verdicts will be provisional or unresolved. % <!-- S2FIX:238-247 ; S2FIX:288-291 -->
\end{itemize}

\paragraph{First-principles calculations.}
\begin{itemize}
\item PBE without a Hubbard correction, spin-unpolarised and neutral cells: the results can show consistency with
donor behaviour, not full ionisation or charge-transition levels. Hybrid functionals enter only as checks, and HSE06
itself underestimates the \InO{} gap. % <!-- REG:260-262 ; REG:585 ; REG:544 -->
\item The absolute gap is anchored to a measurement on a different sputtered IWO device. % <!-- REG:541 -->
\item The machine-learned structures failed the pre-registered force check on the W first shell and are DFT starting
points, not validated amorphous models (claim C037). % <!-- LEDGER:228-233 -->
\item One melt--quench sample per structure (two at \SI{2.0}{nm}); two seeds give diagnostic variability, not a
confidence interval, and a difference smaller than the gap between the two \SI{2.0}{nm} samples is not a result.
Force agreement on a few snapshots would not show that the ensemble represents sputtered films. % <!-- MLV:35-36 ; REG:282 -->
\item Small cells: W images at least \SI{12.5}{\angstrom} apart are only partly controlled, and one excess electron per W
gives about $6\times10^{20}$\cmcu{}, far above the channel carrier density. % <!-- REG:622 ; SLAB:126 -->
\item One state per DFT cell corresponds to about $6\times10^{13}$\cmsq{}, so the fitted fixed charge and interface
trap density of the device model cannot be tested by DFT. % <!-- REG:40 -->
\item Mixed relaxation-smearing families in the W series; a provisional final-SCF width without an electron-affinity
comparison; restart non-reproducibility at production tolerance in one cell (claim C036). % <!-- REG:939-943 ; REG:1025 ; REG:670 ; LEDGER:221-226 -->
\item Deadline effects: relaxations not converged at the cutoff are exploratory; the \SI{1.5}{nm} cells were cut, the
\SI{2.0}{nm} \AlO{} interface cell was stopped and the \SI{2.0}{nm} bare control parked, so H-I is expected to be
unresolved and H-H unresolved at \SI{2.0}{nm}. % <!-- REG:1076 ; DLOG:46 ; REG:823-827 -->
\item H-N is exploratory, and its \SI{2.0}{nm} endpoint is unresolved without a certified C4 polish of both seeds. % <!-- DLOG:39 ; REG:1016-1018 -->
\item The amorphous effective-mass workflow was descoped, so the mass-only mobility prediction is not computed. % <!-- REG:170-171 -->
\item The \HfO{} cell is one seed with a different oxide and bonding; it cannot establish the stack order, and its
oxide density and possible partial crystallisation are assumptions to be checked by the registered flags. % <!-- REG:391-397 ; SLAB:141 ; SLAB:185 -->
\item The sub-nanometre crystalline slab results depend on the termination recipe and are preliminary (claim C021). % <!-- LEDGER:120-124 -->
\item Fixed-geometry hybrid results support electronic robustness only; equilibrium self-trapping would need a
structural qualification, which is not claimed. % <!-- REG:452 -->
\end{itemize}

\paragraph{Procedure.}
\begin{itemize}
\item Partial electronic unblinding of four \SI{2.0}{nm} cells (A36 erratum) and a procedural post-hoc disclosure (A35)
are part of the record (\cref{sec:m-reg-sight}). % <!-- REG:1027 ; REG:950-959 -->
\item TCAD protocol deviations: the S0-02 repair cap and launches while the regression gate was red. % <!-- REG:39 ; S2PREP:82 -->
\item AI reviewers share training priors and can share blind spots; holdout custody, external timestamps and human
judgement are the real safeguards. In this work there is no holdout and the external timestamp is pending. % <!-- CHARTER:84-85 -->
\pdPending{P31}{TEXT PENDING: final wording once the external timestamp (P27) and the charter adoption (P26) are settled}.
\end{itemize}


% =====================================================================================================================
\section*{Data availability}
\pdPending{P37}{TEXT PENDING: data statement; depends on the fabrication group's permission (P01)}

\section*{Code availability}
\pdPending{P38}{TEXT PENDING: code statement (deck generator, scorer, read-out package with freeze tags, structure and
input generators); archive with a DOI}

\section*{Author contributions}
\pdPending{P34}{DECISION PENDING: author list, order and affiliations, agreed with the fabrication group}

\section*{Use of AI tools}
\begin{pdIntent}
Describe the roles of the AI agents as practised (computational execution and record keeping by Claude Code,
independent review by Codex), that every amendment needed a human decision, and that agreement between AI systems was
not treated as evidence. % <!-- CHARTER:15-21 ; CHARTER:53 -->
\end{pdIntent}
\pdPending{P39}{TEXT PENDING: AI-use statement in the wording the journal requires}

\section*{Acknowledgements and competing interests}
\pdPending{P40}{TEXT PENDING: acknowledgements (CERN computing resources, funding) and competing-interest statement}

\bibliographystyle{unsrtnat}
\bibliography{../../../IWO_PHYSICS_CONSTRAINED_MODEL_V1/report_latex_full/references,../../../IWO_PHYSICS_CONSTRAINED_MODEL_V1/report_latex_full/bib/extra_DFT}

\end{document}
